\part{一般理论与沿忠实平坦态射的下降}\label{fga3.i}\thispagestyle{fancy}\begin{original}{}\label{fga3.i-original-citation}\skipsinglepar\par{}A. Grothendieck.
    "Technique de descente et théorèmes d'existence en géométrie algébrique, I: Généralités. Descente par morphismes fidèlement plats".
    \emph{Séminaire Bourbaki} \textbf{12} (1959–60), 第 190 次报告，1959 年 12 月。
    \url{http://www.numdam.org/book-part/SB_1958-1960__5__299_0/}\end{original}\par{}\emph{［编者补注］} 关于下降理论的各种细节，另见［\hyperref[Gro1960b]{Gro1960b}，VI、VII 和 VIII］。\par{}从技术角度看，本文及后续文章可以看作 Hilbert 著名的“定理 90”的变式。将下降方法引入代数几何似乎应归功于 A. Weil，当时称为“基域的下降”。Weil 只考虑有限可分域扩张。随后 P. Cartier 研究了高度为 1 的根式扩张情形。由于缺乏概形语言，尤其是所考虑的环中缺乏幂零元，Cartier 未能表述这两种情形的本质一致性。\par{}目前看来，以下将解释的一般下降技术，必要时结合“形式几何”的基本定理（参见 \hyperref[fga2]{[3]}），是代数几何中大多数存在定理的基础。（\emph{［编者补注］ 如今看来，说下降技术是“代数几何中大多数存在定理的基础”有些过强。对本系列，即“代数几何中的下降技术与存在定理”，前两讲所研究的非射影技术，这在很大程度上成立；但对射影技术，即第 \hyperref[fga3.iv]{IV}、\hyperref[fga3.v]{V} 和 \hyperref[fga3.vi]{VI} 讲，则并非如此。}）还应指出，上述下降技术当然能移植到“解析几何”中，我们可以希望在不太遥远的将来，专家们能够证明第 \hyperref[fga3.ii]{II} 讲将给出的形式几何存在定理的“解析”对应物。无论如何，Kodaira–Spencer 最近的工作已相当清楚地说明，需要更接近概形理论的方法，自然应与超越技术互补，因为他们的方法似乎不适于在“模簇”的奇点邻域内定义并研究这些模簇。\par{}本讲，即第 I 讲，将讨论下降的最初等情形，也就是标题所指的情形。不过，下文 \cref{fga3.i-b.1} 中的 \hyperref[fga3.i-b.1-theorem-1]{定理~\ref*{fga3.i-b.1-theorem-1}}、\hyperref[fga3.i-b.1-theorem-2]{定理~\ref*{fga3.i-b.1-theorem-2}} 与 \hyperref[fga3.i-b.1-theorem-3]{定理~\ref*{fga3.i-b.1-theorem-3}} 的应用已非常多。这里只举其中若干例子，不追求可能的最大一般性。\par{}我们将自由使用概形语言，读者可参阅前引文章及 \hyperref[GR1958]{[2]}。不过须明确指出，本文考虑的预概形不必为 Noether 预概形，态射也不必有限型。例如，若 \(A\) 为完备化为 \(\overline {A}\) 的 Noether 局部环，则需要考虑非 Noether 环 \(\overline {\overline {A}}\otimes _A\overline {A}\)，以及对应于相关环之间的包含关系的仿射概形态射。\chapter{范畴论预备知识}\thispagestyle{fancy}\label{fga3.i-a}\section{纤维范畴、下降资料与 \(\mathcal {F}\) 下降态射}\label{fga3.i-a.1}\subsection{}\label{fga3.i-a.1.a}\begin{definition}{1.1}\setnumber{1.1}\label{fga3.i-a.1-definition-1.1}\skipsinglepar\par{}一个 \emph{以 \(\mathcal {C}\) 为基底的纤维范畴 \(\mathcal {F}\)}（或称 \emph{在 \(\mathcal {C}\) 上} 的纤维范畴）由以下资料组成：\begin{itemize}\item{}一个范畴 \(\mathcal {C}\)；

      \item{}对每个 \(X\in \mathcal {C}\)，一个范畴 \(\mathcal {F}_X\)；

      \item{}对每个 \(\mathcal {C}\) 态射 \(f\colon  X\to  Y\)，一个函子 \(f^*\colon \mathcal {F}_Y\to \mathcal {F}_X\)；对 \(\xi \in \mathcal {F}_Y\)，也将其写作 \[           f^*(\xi ) = \xi  \times _Y X,\]这里把 \(X\) 视为“\(Y\) 上的 \(\mathcal {C}\) 的对象”，即赋有态射 \(f\) 的对象；

      \item{}对任意两个可合成态射 \(X\xrightarrow {f}Y\xrightarrow {g}Z\)，一个函子同构 \[           c_{f,g}\colon  (gf)^* \to  f^*g^*         \]；\end{itemize} 这些资料须满足下列条件：\begin{enumerate}
      
 \item[i.]
        \(\operatorname {id}^*=\operatorname {id}\)
      


      
 \item[ii.]若 \(f\) 或 \(g\) 为恒等同构，则 \(c_{f,g}\) 为恒等同构；


      
 \item[iii.]对任意三个可合成态射 \(X\xrightarrow {f}Y\xrightarrow {g}Z\xrightarrow {h}T\)，利用形如 \(c_{u,v}\) 的同构得到的下图交换：\[           \begin {CD}             (h(gf))^* @= ((hg)f)^*           \\@VVV @VVV           \\(gf)^*h^* @. f^*(hg)^*           \\@VVV @VVV           \\(f^*g^*)h^* @= f^*(g^*h^*)           \end {CD}         \]

    \end{enumerate}\end{definition}\begin{example}{1}\setnumber{1}\label{fga3.i-a.1-example-1}\skipsinglepar\par{}设 \(\mathcal {C}\) 为一切纤维积均存在的范畴。可定义以 \(\mathcal {C}\) 为基底的纤维范畴 \(\mathcal {F}\) 如下：取 \(\mathcal {F}_X\) 为 \(X\) 上的 \(\mathcal {C}\) 的对象范畴，对应于态射 \(f\colon  X\to  Y\) 的函子 \(f^*\colon \mathcal {F}_Y\to \mathcal {F}_X\) 由 \emph{纤维积} \(Z\mapsto  Z\times _Y X\) 定义。\end{example}\begin{example}{2}\setnumber{2}\label{fga3.i-a.1-example-2}\skipsinglepar\par{}设 \(\mathcal {C}\) 为预概形范畴，对 \(X\in \mathcal {C}\)，令 \(\mathcal {F}_X\) 为 \(X\) 上的拟凝聚模层范畴。若 \(f\colon  X\to  Y\) 为预概形态射，则 \(f^*\colon \mathcal {F}_Y\to \mathcal {F}_X\) 为 \emph{模层逆像} 函子。这样便得到 \(\mathcal {C}\) 上的一个纤维范畴。\end{example}\subsection{}\label{fga3.i-a.1.b}\begin{definition}{1.2}\setnumber{1.2}\label{fga3.i-a.1-definition-1.2}\skipsinglepar\par{}若 \(u\) 给出从 \(E\) 到 \(E'\) 中满足 \(v_1(x')=v_2(x')\) 的元素 \(x'\in  E'\) 所成子集的双射，则称集合映射图式 \[       E \xrightarrow {u}       E' \overset {v_1}{\underset {v_2}{\rightrightarrows }}       E''     \] 是 \emph{正合} 的。\end{definition}\begin{definition}{1.3}\setnumber{1.3}\label{fga3.i-a.1-definition-1.3}\skipsinglepar\par{}设 \(\mathcal {F}\) 为以 \(\mathcal {C}\) 为基底的纤维范畴，并考虑 \(\mathcal {C}\) 中满足 \(\alpha \beta _1=\alpha \beta _2\) 的态射图式 \[       S \xleftarrow {\alpha }       S' \overset {\beta _1}{\underset {\beta _2}{\leftleftarrows }}       S''     \]。若对 \(\mathcal {F}_S\) 中的每一对元素 \((\xi ,\eta )\)，集合图式 \begin{eqn}{+}\setnumber{+}\label{fga3.i-a.1-definition-1.3-equation}\skipsinglepar\[         \operatorname {Hom}(\xi ,\eta ) \xrightarrow {\alpha ^*}         \operatorname {Hom}(\alpha ^*(\xi ),\alpha ^*(\eta )) \overset {\beta _1^*}{\underset {\beta _2^*}{\rightrightarrows }}         \operatorname {Hom}(\gamma ^*(\xi ),\gamma ^*(\eta ))       \tag{+}       \]\end{eqn} 都正合，其中 \(\gamma =\alpha \beta _1=\alpha \beta _2\)，则称此图式为 \emph{\(\mathcal {F}\) 正合} 的。\par{}为简便起见，上图利用 \(c_{\beta _i,\alpha }\) 将 \(\beta _i^*\alpha ^*\) 与 \((\alpha \beta _i)^*=\gamma ^*\) 等同。\end{definition}\begin{definition}{1.4}\setnumber{1.4}\label{fga3.i-a.1-definition-1.4}\skipsinglepar\par{}设 \(\mathcal {F}\) 为以 \(\mathcal {C}\) 为基底的纤维范畴，考虑 \(\mathcal {C}\) 中的态射 \(\beta _1,\beta _2\colon  S''\to  S'\)。设 \(\xi '\in \mathcal {F}_{S'}\)。\(\xi '\) 上关于二元组 \((\beta _1,\beta _2)\) 的一个 \emph{粘合资料}，定义为从 \(\beta _1^*(\xi ')\) 到 \(\beta _2^*(\xi ')\) 的一个同构。若 \(\xi ',\eta '\in \mathcal {F}_{S'}\) 均赋有粘合资料，则当下图交换时，称 \(\mathcal {F}_{S'}\) 中的态射 \(u\colon \xi '\to \eta '\) \emph{与粘合资料相容}：\[       \begin {CD}         \beta _1^*(\xi ') @>>> \beta _2^*(\xi ')       \\@VVV @VVV       \\\beta _1^*(\eta ') @>>> \beta _2^*(\eta ').       \end {CD}     \]\end{definition}\par{}由此定义，\(\mathcal {F}_{S'}\) 中赋有关于 \(\beta _1\) 和 \(\beta _2\) 的粘合资料的对象构成一个 \emph{范畴}。若 \(\alpha \colon  S'\to  S\) 是满足 \(\alpha \beta _1=\alpha \beta _2\) 的态射，则对每个 \(\xi \in \mathcal {F}_{S'}\)，\(\mathcal {F}_{S'}\) 的对象 \(\xi '=\alpha ^*(\xi )\) 赋有一个典范粘合资料，因为 \[     \beta _i^*\alpha ^*(\xi )     \simeq  (\alpha \beta _i)^*(\xi )     = \gamma ^*(\xi ),   \]，这里仍置 \(\gamma =\alpha \beta _1=\alpha \beta _2\)；此外，若 \(u\colon \xi \to \eta \) 为 \(\mathcal {F}_S\) 中的态射，则 \[     \alpha ^*(u)\colon  \alpha ^*(\xi ) \to  \alpha ^*(\eta )   \] 是 \(\mathcal {F}_{S'}\) 中与典范粘合资料相容的态射。于是得到从范畴 \(\mathcal {F}_S\) 到 \(\mathcal {F}_{S'}\) 中赋有关于二元组 \((\beta _1,\beta _2)\) 的粘合资料的对象范畴的 \emph{典范函子}。因此，\hyperref[fga3.i-a.1-definition-1.3]{定义~\ref*{fga3.i-a.1-definition-1.3}} 也可重新表述为：当上述函子为 \emph{全忠实} 时，称 \hyperref[fga3.i-a.1-definition-1.3-equation]{式~\ref*{fga3.i-a.1-definition-1.3-equation}} 为 \emph{\(\mathcal {F}\) 正合} 的；即上述函子给出范畴 \(\mathcal {F}_S\) 与 \(\mathcal {F}_{S'}\) 中赋有关于 \((\beta _1,\beta _2)\) 的粘合资料的对象范畴的一个子范畴之间的等价。\begin{definition}{1.5}\setnumber{1.5}\label{fga3.i-a.1-definition-1.5}\skipsinglepar\par{}若赋有此资料的 \(\xi '\) 同构于某个 \(\xi \in \mathcal {F}_S\) 所对应的 \(\alpha ^*(\xi )\)，则称 \(\xi '\in \mathcal {F}_{S'}\) 上的粘合资料关于 \(\alpha \) 是 \emph{有效} 的。\end{definition}\par{}若 \hyperref[fga3.i-a.1-definition-1.3-equation]{式~\ref*{fga3.i-a.1-definition-1.3-equation}} 为 \(\mathcal {F}\) 正合，则 \hyperref[fga3.i-a.1-definition-1.5]{定义~\ref*{fga3.i-a.1-definition-1.5}} 中的对象 \(\xi \) 在唯一同构的意义下确定，并且 \emph{范畴 \(\mathcal {F}_S\) 等价于 \(\mathcal {F}_{S'}\) 中赋有有效粘合资料的对象范畴}。\subsection{}\label{fga3.i-a.1.c}\par{}最重要的情形是 \[     S'' = S' \times _S S',   \]，其中各 \(\beta _i\) 是从 \(S'\times _S S'\) 到其两个因子的投影 \(p_1\) 和 \(p_2\)；这里假设 \(\mathcal {C}\) 中一切纤维积存在。某个 \(\xi '\in \mathcal {F}_{S'}\) 上的粘合资料 \(\varphi \colon  p_1^*(\xi ')\to  p_2^*(\xi ')\) 若有效，立即得到两个必要条件：\begin{enumerate}
    
 \item[i.]\(\Delta ^*(\varphi ) = \operatorname {id}_{\xi '} \)，其中 \(\Delta \colon  S'\to  S'\times _S S'\) 表示对角态射，并将 \(\Delta ^* p_i^*(\xi ')\) 与 \((p_i\Delta )^*(\xi ')=\xi '\) 等同；


    
 \item[ii.]\(p_{31}^*(\varphi ) = p_{32}^*(\varphi )p_{21}^*(\varphi )\)，其中 \(p_{ij}\) 表示从 \(S'\times _S S'\times _S S'\) 到其第 \(i\) 和第 \(j\) 个因子的部分乘积的典范投影。

  \end{enumerate}\begin{definition}{1.6}\setnumber{1.6}\label{fga3.i-a.1-definition-1.6}\skipsinglepar\par{}\(\xi '\in \mathcal {F}_{S'}\) 上关于态射 \(\alpha \colon  S'\to  S\) 的 \emph{下降资料}，定义为 \(\xi '\) 上关于两个典范投影 \(S'\times _S S'\to  S'\) 所成的二元组 \((p_1,p_2)\) 的粘合资料，并要求满足上述条件（i）和（ii）。\end{definition}\begin{definition}{1.7}\setnumber{1.7}\label{fga3.i-a.1-definition-1.7}\skipsinglepar\par{}若态射图式 \[       S \xleftarrow {\alpha }       S' \overset {p_1}{\underset {p_2}{\leftleftarrows }}       S'\times _S S'     \] 为 \(\mathcal {F}\) 正合（\hyperref[fga3.i-a.1-definition-1.3]{定义~\ref*{fga3.i-a.1-definition-1.3}}），则称态射 \(\alpha \colon  S'\to  S\) 为 \emph{\(\mathcal {F}\) 下降态射}；若此外 \(\mathcal {F}_{S'}\) 的任意对象上的每个下降资料（\hyperref[fga3.i-a.1-definition-1.6]{定义~\ref*{fga3.i-a.1-definition-1.6}}）均有效，则称 \(\alpha \) 为 \emph{严格 \(\mathcal {F}\) 下降态射}。后一个严格性条件也可更直观地表述为：“给出 \(\mathcal {F}_S\) 的一个对象，等价于给出 \(\mathcal {F}_{S'}\) 的一个赋有下降资料的对象。”\end{definition}\par{}注意，若 \(\mathcal {F}\) 下降态射（\emph{［编者补注］} 这里不必假设 \(\alpha \) 是 \(\mathcal {F}\) 下降态射）\(\alpha \colon  S'\to  S\) 有一个 \emph{截面} \(s\colon  S\to  S'\)，即满足 \(\alpha  s=\operatorname {id}_S\) 的态射 \(s\)，则它是严格 \(\mathcal {F}\) 下降态射：若 \(\xi '\in \mathcal {F}_{S'}\) 赋有关于 \(\alpha \) 的下降资料，那么“它来自” \(\xi =s^*(\xi ')\)。\subsection{}\label{fga3.i-a.1.d}\par{}可以更直观地表述上述概念：对 \(S\) 上的 \(\mathcal {C}\) 的对象 \(T\)，引入集合 \[     \operatorname {Hom}_S(T,S') = S'(T),   \]，其元素记为 \(t\)、\(t'\) 等。给定对象 \(\xi '\in \mathcal {F}_{S'}\)，每个 \(t\in  S'(T)\) 都对应于 \(\mathcal {F}_T\) 中的对象 \(t^*(\xi ')\)，亦记为 \(\xi '_t\)。于是，\(\xi '\) 上关于 \((p_1,p_2)\) 的粘合资料由如下资料定义：对 \(S\) 上的每个 \(T\) 及每一对点 \(t,t'\in  S'(T)\)，给出同构 \[     \varphi _{t',t}\colon  \xi '_t \to  \xi '_{t'}   \]，并满足关于 \(T\) 的显然的函子性条件。此时 \cref{fga3.i-a.1.c} 的条件（i）和（ii）可写为 \begin{enumerate}
    
 \item[i bis.]对所有 \(T\) 及所有 \(t\in  S'(T)\)，均有 \(\varphi _{t,t}=\operatorname {id}_{\xi '_t}\)；


    
 \item[ii bis.]对所有 \(T\) 及所有 \(t,t',t''\in  S'(T)\)，均有 \(\varphi _{t,t''}=\varphi _{t,t'}\varphi _{t',t''}\)。

  \end{enumerate}。令 \(t=t'=t''\)，可见（ii bis）蕴含 \(\varphi _{t,t}^2=\varphi _{t,t}\)；所以，由于假设 \(\varphi _{t,t}\) 是同构，它蕴含（i bis）。因而（i bis）是（ii bis）的推论，（i）也是（ii）的推论。但若不预先假设各个 \(\varphi _{t,t'}\) 都是同构，即不假设 \(\varphi \colon  p_1^*(\xi ')\to  p_2^*(\xi ')\) 是同构，则（ii bis）未必再蕴含（i bis）；不过，（ii bis）与（i bis）合起来确实蕴含各个 \(\varphi _{t,t'}\) 都是同构，因为此时有 \(\varphi _{t,t'}\varphi _{t',t}=\varphi _{t,t}=\operatorname {id}_{\xi '_t}\)。\section{正合图式、严格满态射、下降态射及例子}\label{fga3.i-a.2}\subsection{}\label{fga3.i-a.2.a}\begin{definition}{2.1}\setnumber{2.1}\label{fga3.i-a.2-definition-2.1}\skipsinglepar\par{}设 \(\mathcal {C}\) 为一个范畴。若对每个 \(Z\in \mathcal {C}\)，相应的集合图式 \[       \operatorname {Hom}(Z,T) \to        \operatorname {Hom}(Z,T') \rightrightarrows        \operatorname {Hom}(Z,T'')     \] 都正合（\hyperref[fga3.i-a.1-definition-1.2]{定义~\ref*{fga3.i-a.1-definition-1.2}}），则称态射图式 \[       T \xrightarrow {\alpha }       T' \overset {\beta _1}{\underset {\beta _2}{\rightrightarrows }}       T''     \] \emph{正合}。此时称 \((T,\alpha )\)，或滥用语言称 \(T\)，为态射对 \((\beta _1,\beta _2)\) 的一个 \emph{核}。\end{definition}\par{}显然，此核在唯一同构的意义下确定。若 \(\mathcal {C}\) 为集合范畴，则上述定义与 \hyperref[fga3.i-a.1-definition-1.2]{定义~\ref*{fga3.i-a.1-definition-1.2}} 相容。对偶地，可以定义 \(\mathcal {C}\) 中态射图式 \[     S \xleftarrow {\alpha }     S' \overset {\beta _1}{\underset {\beta _2}{\leftleftarrows }}     S''   \] 的正合性；此时称 \((S,\alpha )\) 为态射对 \((\beta _1,\beta _2)\) 的一个 \emph{余核}。\begin{definition}{2.2}\setnumber{2.2}\label{fga3.i-a.2-definition-2.2}\skipsinglepar\par{}若态射 \(\alpha \colon  S'\to  S\) 是满态射，并且对每个态射 \(u\colon  S'\to  Z\)，下列必要条件也是 \(u\) 可分解为 \(S'\to  S\to  Z\) 的充分条件，则称其为 \emph{严格满态射}：对每个 \(S''\in \mathcal {C}\) 及每对满足 \(\alpha \beta _1=\alpha \beta _2\) 的态射 \(\beta _1,\beta _2\colon  S''\to  S'\)，均有 \(u\beta _1=u\beta _2\)。\end{definition}\par{}若纤维积 \(S'\times _S S'\) 存在，则这等价于图式 \[     S \xleftarrow {\alpha }     S' \overset {p_1}{\underset {p_2}{\leftleftarrows }}     S'\times _S S'   \] 正合，即 \(S\) 为二元组 \((p_1,p_2)\) 的余核。无论如何，余核态射都是严格满态射。还要指出，若严格满态射也是单态射，则它是同构。将对偶概念 \emph{严格单态射} 的展开留给读者。\par{}为了更精确地说明 \(\mathcal {F}\) 下降态射与严格满态射之间的关系，引入如下定义：\begin{definition}{2.3}\setnumber{2.3}\label{fga3.i-a.2-definition-2.3}\skipsinglepar\par{}若对 \(S\) 上的每个 \(T\)，纤维积 \(T'=S'\times _S T\) 存在且投影 \(T'\to  T\) 是满态射，相应地是严格满态射，则称态射 \(\alpha \colon  S'\to  S\) 为 \emph{普遍满态射}，相应地称为 \emph{普遍严格满态射}。\end{definition}\par{}在性质很好的范畴中，例如集合范畴、拓扑空间上的层范畴、abelian 范畴等等，上述四种“满性”概念都一致；但在预概形范畴，或给定非空预概形 \(S\) 上的预概形范畴中，它们各不相同，即使只考虑在 \(S\) 上有限的 \(S\) 概形也是如此。\begin{definition}{2.4}\setnumber{2.4}\label{fga3.i-a.2-definition-2.4}\skipsinglepar\par{}若态射 \(\alpha \colon  S'\to  S\) 是 \(\mathcal {F}\) 下降态射，相应地是严格 \(\mathcal {F}\) 下降态射（参见 \hyperref[fga3.i-a.1-definition-1.7]{定义~\ref*{fga3.i-a.1-definition-1.7}}），则称其为 \emph{下降态射}，相应地称为 \emph{严格下降态射}；其中 \(\mathcal {F}\) 是 \(\mathcal {C}\) 上的纤维范畴，其纤维由 \(\mathcal {C}\) 的对象上的 \(\mathcal {C}\) 的对象组成（参见 \hyperref[fga3.i-a.1-example-1]{例~\ref*{fga3.i-a.1-example-1}}）。\end{definition}\begin{proposition}{2.1}\setnumber{2.1}\label{fga3.i-a.2-proposition-2.1}\skipsinglepar\par{}若 \(\mathcal {C}\) 中一切有限乘积和有限纤维积均存在，则 \(\mathcal {C}\) 中的下降态射与 \(\mathcal {C}\) 中的普遍严格满态射一致。\end{proposition}\subsection{}\label{fga3.i-a.2.b}\begin{example}{}\skipsinglepar\par{}设 \(\mathcal {C}\) 为预概形范畴，\(S\in \mathcal {C}\)，且 \(S'\)、\(S''\) 是 \(S\) 上的 \emph{有限} 预概形；即它们对应于 \(S\) 上的代数层 \(\mathscr {A}'\)、\(\mathscr {A}''\)，作为模层拟凝聚且有限型，若 \(S\) 局部 Noether 则为凝聚层。设 \(\alpha \colon  S'\to  S\) 为 \(S'\) 的结构态射，\(\beta _1\) 和 \(\beta _2\) 为从 \(S''\) 到 \(S'\) 的 \(S\) 态射，由代数同态 \(\mathscr {A}'\to \mathscr {A}''\) 定义，后者也记为 \(\beta _1\) 和 \(\beta _2\)。利用有限态射是闭映射这一事实，即第一 Cohen–Seidenberg 定理，容易证明 \(\mathcal {C}\) 中的图式 \begin{eqn}{+}\setnumber{+}\label{fga3.i-a.2.b-equation}\skipsinglepar\[         S \xleftarrow {\alpha }         S' \overset {\beta _1}{\underset {\beta _2}{\leftleftarrows }}         S''       \tag{+}       \]\end{eqn} 正合，当且仅当 \(S\) 上的层图式 \[       \mathscr {O}_S = \mathscr {A} \xrightarrow {\alpha }       \mathscr {A}' \overset {\beta _1}{\underset {\beta _2}{\rightrightarrows }}       \mathscr {A}''     \] 正合。特别地，若 \(\alpha \colon  S'\to  S\) 为对应于 \(S\) 上代数层 \(\mathscr {A}'\) 的有限态射，则 \(\alpha \) 是严格满态射，当且仅当层图式 \[       \mathscr {O}_S = \mathscr {A} \to        \mathscr {A}' \overset {p_1}{\underset {p_2}{\rightrightarrows }}       \mathscr {A}'\otimes _{\mathscr {A}}\mathscr {A}'     \] 正合；它是满态射，当且仅当 \(\mathscr {A}\to \mathscr {A}'\) 为单射。若 \(S\) 是环为 \(A\) 的仿射概形，则 \(S'\) 是环为 \(A'\) 的仿射概形，\(A'\) 在 \(A\) 上有限；所以 \(S'\to  S\) 是严格满态射，当且仅当 \(A\to  A'\) 给出从 \(A\) 到 \(A'\) 中满足 \[       1_{A'}\otimes _A x' - x'\otimes _A 1_{A'}       = 0     \] 的元素 \(x'\in  A'\) 所成子环的同构；它是满态射，当且仅当 \(A\to  A'\) 为单射。如前所述，即使 \(S\) 是 Artin 局部环的概形，一个为满态射的有限态射 \(S'\to  S\) 也未必为严格满态射。不过，可以证明，\emph{若 \(S\) 为 Noether 预概形，则每个为满态射的有限态射 \(S'\to  S\) 都是有限个严格满态射的复合}，其中各态射也有限。这也表明，两个严格满态射的复合未必是严格满态射。\end{example}\subsection{}\label{fga3.i-a.2.c}\par{}若 \hyperref[fga3.i-a.2.b-equation]{式~\ref*{fga3.i-a.2.b-equation}} 是有限态射的正合图式，则对每个预概形的 \emph{平坦} 态射 \(T\to  S\)，由基底变换 \(T\to  S\) 从 \hyperref[fga3.i-a.2.b-equation]{式~\ref*{fga3.i-a.2.b-equation}} 诱导的图式仍正合。因此，若 \(X\)、\(Y\) 是 \(S\) 预概形，且 \(X\) 在 \(S\) 上 \emph{平坦}，则如下映射图式正合：\[     \operatorname {Hom}_S(X,Y) \to      \operatorname {Hom}_{S'}(X',Y') \rightrightarrows      \operatorname {Hom}_{S''}(X'',Y'').   \] 其中 \(X'\)、\(Y'\) 是 \(X\)、\(Y\) 在 \(S'\) 上的逆像，而 \(X''\)、\(Y''\) 是它们在 \(S''\) 上的逆像。特别地，若 \(\mathcal {F}\) 表示预概形范畴 \(\mathcal {C}\) 上的纤维范畴，且对 \(X\in \mathcal {C}\)，\(\mathcal {F}_X\) 为 \emph{平坦} \(X\) 预概形的范畴，则图式 \hyperref[fga3.i-a.2.b-equation]{式~\ref*{fga3.i-a.2.b-equation}} 是 \(\mathcal {F}\) 正合的。若不加平坦性假设，此结论不成立；特别地，有限严格满态射未必为下降态射。类似地，若 \(\mathcal {F}\) 表示使 \(\mathcal {F}_X\) 为预概形 \(X\) 上 \emph{平坦} 拟凝聚层范畴的纤维范畴，则 \hyperref[fga3.i-a.2.b-equation]{式~\ref*{fga3.i-a.2.b-equation}} 是 \(\mathcal {F}\) 正合的；这里平坦性假设同样是必要的。在这两种情形下，\(S'\) 上的平坦对象的粘合资料是否 \emph{有效}，特别是在 \(S''=S'\times _S S'\) 时下降资料是否有效，都是精细的问题；本系列文章的主要目标之一就是在许多特殊情形下回答它。报告人不知道：对每个有限严格满态射 \(S'\to  S\)，\(S'\) 上平坦拟凝聚层的每个下降资料是否都有效；即使假设 \(S\) 是 Artin 局部环的谱，并只考虑秩为 \(1\) 的局部自由层，也不知道。更一般地，设 \(A\) 为环，\(A'\) 为 \(A\) 代数，均交换，使图式 \[     A \to      A' \rightrightarrows      A'\otimes _A A'   \] 正合；这等价于说，相应的 \(A'\) 与 \(A\) 的谱之间的态射 \(S'\to  S\) 是 \(\mathcal {F}\) 下降态射，其中 \(\mathcal {F}\) 为平坦拟凝聚层的纤维范畴。设 \(M'\) 为赋有向 \(A\) 的下降资料的平坦 \(A'\) 模，即赋有 \((A'\otimes _A A')\) 模同构 \[     \varphi \colon  M'\otimes _A A'     \xrightarrow {\sim } A'\otimes _A M'   \]，且满足 \cref{fga3.i-a.1.c} 的条件（i）和（ii）；将它们写成模的形式留给读者。此资料相对于平坦拟凝聚层的纤维范畴是否有效？设 \(M\) 为 \(M'\) 中满足 \[     \varphi (x'\otimes _A 1_{A'})     = 1_{A'}\otimes _A x',   \] 的元素 \(x'\in  M'\) 所成子集，它是 \(M'\) 的 \(A\) 子模。典范单射 \(M\to  M'\) 定义 \(A'\) 模同态 \(M\otimes _A A'\to  M'\)。\emph{于是 \(\varphi \) 的有效性等价于如下事实：\(M\) 为平坦 \(A\) 模，且上述同态为同构。}\begin{remark}{}\label{fga3.i-a.2.c-remark}\skipsinglepar\par{}上述讨论并未对图式 \hyperref[fga3.i-a.2.b-equation]{式~\ref*{fga3.i-a.2.b-equation}} 的态射作任何平坦性假设，因此，为建立下降技术，必须对所考虑的 \(S\) 和 \(S'\) 上的对象作平坦性假设。在 \cref{fga3.i-b.2} 中，将对 \(\alpha \colon  S'\to  S\) 作平坦性假设，从而可对 \(S\) 和 \(S'\) 上不再要求平坦的对象建立下降技术。无论如何，总有某个地方涉及平坦性假设。这是平坦性概念在代数几何中如此重要的主要原因之一；只考虑以 \emph{域} 为基底时，这种作用不会显现，因为域上一切对象实际上都平坦！\end{remark}\section{对 étale 态射的应用}\label{fga3.i-a.3}\par{}设 \(A\) 为局部环，\(B\) 为 \(A\) 上的局部代数，其极大理想诱导 \(A\) 的极大理想。若满足以下条件，则称 \(B\) 在 \(A\) 上 \emph{étale}，而不用其他地方所说的“非分歧”：\begin{enumerate}
    
 \item[i.]\(B\) 在 \(A\) 上平坦；

    
 \item[ii.]\(B/\mathfrak {m}B\) 为 \(A/\mathfrak {m}=k\) 的有限可分扩张，其中 \(\mathfrak {m}\) 表示 \(A\) 的极大理想。

  \end{enumerate} 若 \(A\)、\(B\) 均为 Noether 环，且 \(k\) 代数闭，则这等价于延拓 \(A\to  B\) 的完备化同态 \(\overline {A}\to \overline {B}\) 是同构。若 \(\mathscr {O}_x\) 在 \(\mathscr {O}_{f(x)}\) 上 étale，则称有限型态射 \(f\colon  T\to  S\) \emph{在 \(x\in  T\) 处 étale}，或称 \(T\) \emph{在 \(x\) 处相对于 \(S\) 为 étale}；若 \(f\) 在每个 \(x\in  T\) 处 étale，则称 \(f\) 为 \emph{étale} 态射，或 \emph{étale 态射}，也称 \(T\) \emph{在 \(S\) 上 étale}。注意，若 \(S\) 局部 Noether，则 \(T\) 中使 \(f\) 为 étale 的点构成开集，Zariski“主定理”使我们能够精确描述 \(T/S\) 在这种点的邻域内的结构，给出一个众所周知类型的方程。\par{}若 \(S\) 为复数域上的有限型概形，则对应于一个解析空间 \(\overline {S}\)，其意义与 Serre \hyperref[Ser1956]{[5]} 相同，只是 \(\overline {S}\) 的结构层可以有幂零元，而这并不实质性改变 \hyperref[Ser1956]{[5]} 的内容。容易看出，\(f\) 是 étale 态射，当且仅当 \(\overline {f}\colon \overline {T}\to \overline {S}\) 是 étale 态射，即 \(\overline {T}\) 的每一点都有一个邻域，使 \(\overline {f}\) 在该邻域上诱导到 \(\overline {S}\) 的开子集的同构。特别地，\(S\) 的每个 \emph{étale 覆盖} \(T\)，即每个有限 étale 态射 \(f\colon  T\to  S\)，都对应于 \(\overline {S}\) 的一个 étale 覆盖 \(\overline {T}\)；它连通，当且仅当 \(T\) 连通 \hyperref[Ser1956]{[5]}。还容易看出，若 \(T\)、\(T'\) 为 \(S\) 上的 étale 概形，则自然映射 \[     \operatorname {Hom}_S(T,T') \to  \operatorname {Hom}_{\overline {S}}(\overline {T},\overline {T}')   \] 是双射；即从 \(S\) 上的 étale 概形范畴到 \(\overline {S}\) 上的 étale 解析空间范畴的函子 \(T\mapsto \overline {T}\) 是“全忠实”的，因此给出前一个范畴与后一个范畴的一个子范畴之间的等价。Grauert–Remmert 的定理 \hyperref[GR1958]{[2]} 蕴含：若 \(S\) 正规，则由此得到 \(S\) 的 \emph{étale 覆盖} 范畴与 \(\overline {S}\) 的 \emph{有限} étale 覆盖范畴之间的等价；即 \(\overline {S}\) 的每个 étale 覆盖 \(\mathscr {C}\) 均与某个 \(\overline {T}\) 在 \(\overline {S}\) 上同构，其中 \(T\) 为 \(S\) 的 étale 覆盖。我们将证明 \emph{不对 \(S\) 作任何正规性假设，Grauert–Remmert 定理仍成立}。先设 \(S'\to  S\) 为有限严格满态射，并假设定理对 \(S'\) 已成立；要证明它对 \(S\) 成立。设 \(\mathscr {C}\) 为 \(\overline {S}\) 的 étale 覆盖，考虑它在 \(S'\) 上的逆像 \(\mathscr {C}'\)；后者对应于 \(S'\) 上的一个凝聚解析代数层 \(\mathfrak {A}'\)，该层为 \(\overline {S}\) 上定义 \(\mathscr {C}\) 的代数层 \(\mathfrak {A}\) 的逆像。由关于 \(S'\) 的假设，\(\mathscr {C}'\) 来自 \(S'\) 的一个 étale 覆盖 \(T'\)，即 \(\mathfrak {A}'\) 来自 \(S'\) 上的凝聚代数层 \(\mathscr {A}'\)。此外，\(\mathfrak {A}'\) 赋有关于 \(\overline {S}'\to \overline {S}\) 的典范下降资料，即其在 \(\overline {S}'\times _{\overline {S}}\overline {S}'=\overline {S'\times _SS'}\) 上的两个逆像之间的同构，满足条件（i）和（ii）；由前述讨论，此同构来自相应代数层之间的同构，即来自 \(\mathscr {A}'\) 上关于 \(S'\to  S\) 的下降资料。容易证明后者有效，因为它在 \(\mathfrak {A}'\) 上有效，而上一节已明确描述的下降资料的有效性，可以在所涉及的模的 \emph{完备化} 上局部检验。由此得到 \(S\) 上的凝聚代数层 \(\mathscr {A}\)，定义了 \(S\) 的一个覆盖 \(T\)，这就是所求覆盖。显然，当 \(S'\to  S\) 仅为有限多个有限严格满态射的复合，即由 \cref{fga3.i-a.2} 的分解结果，仅为任意有限满态射时，上述结果仍成立。因此，若 \(S\) 为 \emph{既约} 概形，即 \(\mathscr {O}_S\) 没有幂零元，则引入其正规化 \(S'\)，便知 Grauert–Remmert 定理仍成立。容易过渡到一般情形。\par{}一个完全类似的论证，同样利用有限严格满态射的分解结果以及下降资料有效性的“形式”性质，能证明如下结果：设 \(S\) 为局部 Noether 预概形，\(S'\to  S\) 为有限、满射且根式的态射；等价地，它是一个有限型态射，使得对 \(S\) 上的每个 \(T\)，态射 \(T'=S'\times _S T\to  T\) 都是同胚，也可说 \(S'\to  S\) 是 \emph{普遍同胚}。对 \(S\) 上的每个 étale 预概形 \(T\)，考虑其逆像 \(T'=T\times _S S'\)，它在 \(S'\) 上 étale。\emph{于是函子 \(T\mapsto  T'\) 是 \(S\) 上的 étale 预概形 \(T\) 的范畴与 \(S'\) 上的 étale 预概形 \(T'\) 的范畴之间的等价。} 这里使用了：对 \(S\) 上的 étale 预概形 \(T_1\) 和 \(T_2\)，映射 \[     \operatorname {Hom}_S(T_1,T_2) \to  \operatorname {Hom}_{S'}(T'_1,T'_2)   \] 为双射；这可毫无困难地直接证明。还使用了：若 \(S'=(S,\mathscr {O}_S/\mathscr {J})\)，其中 \(\mathscr {J}\) 为 \(\mathscr {O}_S\) 的幂零凝聚理想层，则所述定理成立，参见［\hyperref[Mur1958]{4}，引理 6］。还要注意，这里不假设所涉及的 \(T\) 在 \(S\) 上有限。特别地，此结果蕴含态射 \(S'\to  S\) 诱导 \(S'\) 的基本群与 \(S\) 的基本群之间的同构，即“\emph{预概形基本群的拓扑不变性}”。\section{与一次上同调的关系}\label{fga3.i-a.4}\subsection{}\label{fga3.i-a.4.a}\par{}设 \(\mathcal {C}\) 为任意两个对象的乘积总存在的范畴，且 \(T\in \mathcal {C}\)。对每个有限集 \(I\neq \varnothing \)，可以考虑 \(T^I\)，从而得到从非空有限集范畴到范畴 \(\mathcal {C}\) 的协变函子，即可称为 \(\mathcal {C}\) 的一个 \emph{单纯对象}，记为 \(K_T\)。此对象协变地依赖于 \(T\)；此外，\emph{若 \(u\) 和 \(v\) 是态射 \(T\to  T'\)，则相应态射 \(K_T\to  K_{T'}\) 同伦}。若 \(\operatorname {Hom}(T,T')\neq \varnothing \)，则称 \(T\) \emph{支配} \(T'\)，这在 \(\mathcal {C}\) 上给出一个向上有向的预序。由前述讨论，若 \(T\) 支配 \(T'\)，则存在单纯对象同态 \(K_T\to  K_{T'}\) 的一个典范同伦类；特别地，若 \(K_T\)、\(K_{T'}\) 所对应的 \(T\) 与 \(T'\) 相互支配，则 \(K_T\) 与 \(K_{T'}\) 同伦等价。现设 \(F\) 为从 \(\mathcal {C}\) 到 \emph{abelian} 范畴 \(\mathcal {C}'\) 的函子，明确起见，设其反变。于是 \[     C^\bullet (T,F) = F(K_T)   \] 是 \(\mathcal {C}'\) 的余单纯对象，从而以熟知的方式定义 \(\mathcal {C}'\) 中的上链复形，可取其上同调：\[     \operatorname {H}^\bullet (T,F)     = \operatorname {H}^\bullet (C^\bullet (T,F))     = \operatorname {H}^\bullet (F(K_T))   \] 若可能混淆，可在 \(\operatorname {H}^\bullet \) 上添下标“\(\mathcal {C}\)”。这是关于 \(F\) 的上同调函子；当 \(T\) 变化时的变性，由先前关于 \(K_T\) 的讨论得出。更确切地说，固定 \(F\)，让 \(T\) 在以支配关系预序的 \(\mathcal {C}\) 中变化，则各个 \(\operatorname {H}^\bullet (T,F)\) 构成 \(\mathcal {C}'\) 中分次对象的归纳系；特别地，若 \(T\)、\(T'\) 相互支配，则 \(\operatorname {H}^\bullet (T,F)\) 与 \(\operatorname {H}^\bullet (T',F)\) 典范同构。\par{}假设 \(\mathcal {C}\) 中一切纤维积均存在。固定 \(S\in \mathcal {C}\)，即可将上述讨论用于 \(S\) 上的 \(\mathcal {C}\) 的对象范畴 \(\mathcal {C}_S\)；若要明确指出是在范畴 \(\mathcal {C}_S\) 中工作，就用 \(C^\bullet (T/S,F)\) 与 \(\operatorname {H}^\bullet (T/S,F)\) 代替 \(C^\bullet (T,F)\) 与 \(\operatorname {H}^\bullet (T,F)\)。于是 \(C^\bullet (T/S,F)\) 是 \(\mathcal {C}'\) 中的上链复形，其 \(n\) 次项为 \(F(T\times _S T\times _S\ldots \times _S T)\)，其中共有 \(n+1\) 个因子 \(T\)。\par{}注意，照常可以不假设范畴 \(\mathcal {C}'\) 为 abelian 范畴就定义 \(\operatorname {H}^0(T/S,F)\)：它是对应于两个投影 \(p_1,p_2\colon  T\times _S T\to  T\) 的态射对 \((F(p_1),F(p_2))\)、\[     F(T) \to  F(T\times _S T)   \] 的核（\hyperref[fga3.i-a.2-definition-2.1]{定义~\ref*{fga3.i-a.2-definition-2.1}}），只要此核存在。特别地，有自然态射 \[     F(S) \to  \operatorname {H}^0(T/S,F)   \]，称为 \emph{增广}；在良好情形下它是同构，特别是当 \(T\to  S\) 为严格满态射且 \(F\)“左正合”时如此。类似地，若 \(F\) 取值于某范畴 \(\mathcal {C}''\) 中的群范畴，也可定义 \(\operatorname {H}^1(T/S,F)\)；当 \(\mathcal {C}''\) 为集合范畴，即 \(F\) 取值于不必交换的群范畴时，\(\operatorname {H}^1(T,F)\) 是 \(C^1(T/S,F) = F(T\times _S T)\) 中满足 \[     F(p_{31})(g)     = F(p_{32})(g) F(p_{21})(g)   \] 的元素 \(g\) 所成子群 \(Z^1(T/S,F)\)，对算子群 \(F(T)\) 的作用取商。该群通过 \[     \rho (g')\cdot  g     = F(p_2)(g') g F(p_1)(g')^{-1}.   \] 作用于 \(C^1(T/S,F)\)，因而特别作用于子集 \(Z^1(T/S,F)\)。\subsection{}\label{fga3.i-a.4.b}\par{}例如，设 \(\mathcal {F}\) 为以 \(\mathcal {C}\) 为基底的纤维范畴，\(\xi ,\eta \in \mathcal {F}_S\)；对 \(S\) 上的所有 \(S'\)，置 \[     F_{\xi ,\eta }(S')     = \operatorname {Hom}(\xi \times _S S', \eta \times _S S').   \]。于是 \(F_{\xi ,\eta }\) 为从 \(\mathcal {C}_S\) 到集合范畴的反变函子。在此设定下，\emph{断言增广态射} \[     F_{\xi ,\eta }(S) \to  \operatorname {H}^0(S'/S,F_{\xi ,\eta })   \] \emph{对每对元素 \(\xi ,\eta \in \mathcal {F}_S\) 都是同构，就等价于断言 \(\alpha \colon  S'\to  S\) 是 \(\mathcal {F}\) 下降态射}（\hyperref[fga3.i-a.1-definition-1.7]{定义~\ref*{fga3.i-a.1-definition-1.7}}）。\subsection{}\label{fga3.i-a.4.c}\par{}类似地，对 \(\xi \in \mathcal {F}_S\)，对 \(S\) 上的 \(\mathcal {C}\) 的任意对象 \(S'\)，用 \[     G_\xi (S')     = \operatorname {Aut}(\xi \times _S S'),   \] 定义从 \(\mathcal {C}_S\) 到群范畴的反变函子 \(G_\xi \)。在此设定下，断言 \emph{\(Z^1(S'/S,G)\) 典范地等同于 \(\xi '=\xi \times _S S'\) 上关于 \(S'\to  S\) 的下降资料的集合}（\hyperref[fga3.i-a.1-definition-1.6]{定义~\ref*{fga3.i-a.1-definition-1.6}}），并且 \emph{\(\operatorname {H}^1(S'/S,G)\) 可等同于如下对象的同构类集：\(\mathcal {F}_{S'}\) 的对象赋有关于 \(\alpha \colon  S'\to  S\) 的下降资料，且作为 \(\mathcal {F}_{S'}\) 的对象同构于 \(\xi '=\xi \times _S S'\)}。于是，\emph{若 \(\alpha \colon  S'\to  S\) 为 \(\mathcal {F}\) 下降态射}（参见 \cref{fga3.i-a.4.b}），\emph{则 \(\operatorname {H}^1(S'/S,G)\) 包含如下对象的同构类集为子集：\(\mathcal {F}_S\) 的对象 \(\eta \) 满足 \(\eta \times _S S'\) 在 \(\mathcal {F}_{S'}\) 中同构于 \(\xi \times _S S'\)}；此外，\emph{此包含为恒等，当且仅当 \(\xi '=\xi \times _S S'\) 上关于 \(\alpha \colon  S'\to  S\) 的每个下降资料都有效}。特别地，若 \(\alpha \colon  S'\to  S\) 是严格 \(S\) 下降态射，就会如此。\begin{remark}{}\label{fga3.i-a.4.c-remark}\skipsinglepar\par{}形如 \(C^\bullet (T/S,F)\) 的上链复形包含大多数已知的标准复形为特例，例如 Čech 上同调、群上同调的复形等，并在代数几何中，特别是在预概形的“Weil 上同调”中发挥重要作用。\end{remark}\subsection{}\label{fga3.i-a.4.d}\begin{example}{1}\setnumber{1}\label{fga3.i-a.4-example-1}\skipsinglepar\par{}设 \(S'\) 为 \(S\in \mathcal {C}\) 上的对象，\(\Gamma \) 为 \(S'\) 的一个自同构群，使得 \(S'\)“在 \(S\) 上形式地为 \(\Gamma \) 主对象”，即自然态射 \[       \Gamma \times  S' \to  S'\times _S S'     \] 是同构，其中 \(\Gamma \times  S'\) 表示 \(\Gamma \) 个 \(S'\) 的直和；假设 \(\mathcal {C}\) 中所需直和均存在。设 \(F\) 为从 \(\mathcal {C}\) 到 abelian 群范畴的反变函子。则 \emph{\(C^\bullet (S'/S,F)\) 典范同构于标准齐次上链的单纯群 \(C^\bullet (\Gamma ,F(S'))\)，所以 \(\operatorname {H}^\bullet (S'/S,F)\) 典范同构于 \(\operatorname {H}^\bullet (\Gamma ,F(S'))\)}。\end{example}\subsection{}\label{fga3.i-a.4.e}\begin{example}{2}\setnumber{2}\label{fga3.i-a.4-example-2}\skipsinglepar\par{}设 \(\mathcal {C}\) 为预概形范畴。以 \(\operatorname {G_a}\) 表示从 \(\mathcal {C}\) 到 abelian 群范畴的反变函子，由 \[       \operatorname {G_a}(X)       = \operatorname {H}^0(X,\mathscr {O}_X).     \] 定义，其记号意为“加法群”。类似地，以 \[       \operatorname {G_m}(X)       = \operatorname {H}^0(X,\mathscr {O}_X)^\times      \] 定义函子 \(\operatorname {G_m}\)，意为“乘法群”，即环 \(\operatorname {H}^0(X,\mathscr {O}_X)\) 的可逆元群；更一般地，以 \[       \operatorname {GL}(n)(X)       = \operatorname {GL}(n,\operatorname {H}^0(X,\mathscr {O}_X)),     \] 定义函子 \(\operatorname {GL}(n)\)，意为“\(n\) 阶线性群”，这是从 \(\mathcal {C}\) 到群范畴的函子，当 \(n>1\) 时群不必交换；取 \(n=1\) 就恢复 \(\operatorname {G_m}\)。还可以将 \(\operatorname {GL}(n)\) 看作自同构函子（参见 \cref{fga3.i-a.4.c}）：考虑以 \(\mathcal {C}\) 为基底的纤维范畴 \(\mathcal {F}\)，其中对 \(X\in \mathcal {C}\)，\(\mathcal {F}_X\) 为 \(X\) 上的局部自由层范畴，因为此时 \(\operatorname {GL}(n)(X)=\operatorname {Aut}_{\mathcal {F}_X}(\mathscr {O}_X^n)\)。由 \cref{fga3.i-a.4.b}，若 \(\alpha \colon  S'\to  S\) 为 \(\mathcal {F}\) 下降态射（参见 \cref{fga3.i-a.2.c}），则 \(\operatorname {H}^1(S'/S,\operatorname {GL}(n))\) 包含 \(S\) 上逆像到 \(S'\) 后同构于 \(\mathscr {O}_{S'}^n\) 的局部自由层的同构类集；此包含为等号，当且仅当 \(\mathscr {O}_{S'}^n\) 上关于 \(\alpha \colon  S'\to  S\) 的每个下降资料均有效。若 \(S\) 为局部环的谱，则这蕴含 \(\operatorname {H}^1(S'/S,\operatorname {GL}(n))=(e)\)，因为此时 \(S\) 上的每个局部自由层均平凡。\par{}注意，关于态射 \(\alpha \colon  S'\to  S\) 的下列条件等价：\begin{enumerate}
      
 \item[i.]增广同态 \(\operatorname {H}^0(S,\mathscr {O}_S) = \operatorname {G_a}(S)\to \operatorname {H}^0(S'/S,\operatorname {G_a})\) 是同构；

      
 \item[ii.]\(\alpha \colon  S'\to  S\) 是 \(\mathcal {F}\) 下降态射，其中 \(\mathcal {F}\) 是上述 \(\mathcal {C}\) 上的纤维范畴；

      
 \item[iii.]若 \(S'\to S\) 有限，则上述条件还等价于：\(\alpha \colon  S'\to  S\) 是严格满态射（参见 \cref{fga3.i-a.2.c}）。

    \end{enumerate}\par{}现在假设 \(S=\operatorname {Spec}(A)\) 和 \(S'=\operatorname {Spec}(A')\)；则 \[       C^n(S'/S,\operatorname {G_a})       = C^n(A'/A,\operatorname {G_a})       = \underbrace {A'\otimes _A A'\otimes _A\ldots \otimes _A A'}_{n+1\text { 个因子 }A'}     \]，其上边缘算子 \(C^n(A'/A,\operatorname {G_a})\to  C^{n+1}(A'/A,\operatorname {G_a})\) 是面算子 \[       \partial _i(x_0\otimes  x_1\otimes \ldots \otimes  x_n)       = x_0\otimes \ldots \otimes  x_{i-1}\otimes 1_{A'}\otimes  x_i\otimes \ldots \otimes  x_n.     \] 的交错和。类似地，\(C^n(S'/S,\operatorname {G_m})=C^n(A'/A,\operatorname {G_m})\) 可等同于 \((\bigotimes _A^{n+1}A')^\times \)，而 \(C^\bullet (A'/A,\operatorname {G_m})\) 的单纯运算由 \(C^\bullet (S'/S,\operatorname {G_a})\) 中的运算诱导。也可同样明确地写出 \(C^\bullet (A'/A,\operatorname {GL}(n))\) 的单纯运算。\emph{在报告人所知道的一切情形下，对 \(i>0\) 有 \(\operatorname {H}^i(A'/A,\operatorname {G_a})=0\)；若 \(A\) 为局部环，则 \(\operatorname {H}^1(A'/A,\operatorname {G_m})=0\)，更一般地，\(\operatorname {H}^1(A'/A,\operatorname {GL}(n))=(e)\)} 若 \(S'\to  S\) 为 \(\mathcal {F}\) 下降态射，即图式 \(A\to  A'\rightrightarrows  A'\otimes _A A'\) 正合，则可与 \cref{fga3.i-a.2.c} 比较。注意，\emph{Hilbert 的“定理 90”恰好就是说：若 \(A\) 为域，\(A'\) 为 \(A\) 的有限 Galois 扩张，则 \(\operatorname {H}^1(S'/S,\operatorname {G_m})=0\)}（参见 \hyperref[fga3.i-a.4-example-1]{例~\ref*{fga3.i-a.4-example-1}}），\emph{也可表述为，在所考虑的情形下，\(S'\to  S\) 是秩为 \(1\) 的局部自由层的纤维范畴的严格下降态射}。最后这种表述最适于推广 Hilbert 定理：既改变对态射 \(S'\to  S\) 的假设，也改变所涉及的拟凝聚层。\par{}最后指出，对极大理想为 \(\mathfrak {m}\) 的局部 \emph{Artin} 环 \(A\) 及 \(A\) 代数 \(A'\)，下列性质等价；这里对每个整数 \(k>0\)，将环 \(A/\mathfrak {m}^{k+1}\)、相应地 \(A'/\mathfrak {m}^{k+1}A'\) 记为 \(A_k\)、相应地 \(A'_k\)：\begin{enumerate}
      
 \item[i.]对所有 \(k\)，有 \(\operatorname {H}^1(A'_k/A_k,\operatorname {G_a})=0\)；

      
 \item[ii.]对所有 \(k\)，有 \(\operatorname {H}^1(A'_k/A_k,\operatorname {G_m})=0\)；

      
 \item[iii.]对所有 \(k\) 和所有 \(n\)，有 \(\operatorname {H}^1(A'_k/A_k,\operatorname {GL}(n))=(e)\)。

    \end{enumerate} 若 \(S'\to  S\) 为严格满态射，则上述条件蕴含它对 \(A'\) 上不必有限型的自由模是 \emph{严格} 下降态射。\end{example}\begin{remark}{}\label{fga3.i-a.4.e-remark}\skipsinglepar\par{}当 \(S\)、相应地 \(S'\) 为域 \(A\)、相应地 \(A'\) 的谱时，群 \(\operatorname {H}^i(S'/S,\operatorname {G_m})\) 的定义归功于 Amitsur。群 \(\operatorname {H}^2(S'/S,\operatorname {G_m})\) 作为 Brauer 群的“整体”变式尤其有趣，参见［\hyperref[GD1960]{1}，VII］。\end{remark}\chapter{沿忠实平坦态射的下降}\thispagestyle{fancy}\label{fga3.i-b}\section{下降定理的陈述}\label{fga3.i-b.1}\begin{definition}{1.1}\setnumber{1.1}\label{fga3.i-b.1-definition-1.1}\skipsinglepar\par{}若对所有 \(x'\in  S'\)，\(\mathscr {O}_{x'}\) 都是环 \(\mathscr {O}_{\alpha (x')}\) 上的平坦模，即 \(\mathscr {O}_{x'}\otimes _{\mathscr {O}_{\alpha (x')}}M\) 是关于 \(\mathscr {O}_{\alpha (x')}\) 模 \(M\) 的正合函子，则称预概形态射 \(\alpha \colon  S'\to  S\) 为 \emph{平坦} 的。若一个态射平坦且满射，则称其为 \emph{忠实平坦} 的。\end{definition}\par{}例如，若 \(S=\operatorname {Spec}(A)\)、\(S'=\operatorname {Spec}(A')\)，则 \(S'\) 在 \(S\) 上平坦，当且仅当 \(A'\) 为平坦 \(A\) 模；\(S'\) 在 \(S\) 上忠实平坦，当且仅当 \(A'\) 为忠实平坦 \(A\) 模，即 \(A'\otimes _A M\) 是关于 \(A\) 模 \(M\) 的 \emph{正合} 且 \emph{忠实} 的函子。这也就是说，按 Serre \hyperref[Ser1956]{[5]} 的术语，二元组 \((A,A')\) 平坦。若 \(S'\) 在 \(S\) 上忠实平坦，则 \(S\) 上拟凝聚层的逆像函子正合且忠实；换言之，\(S\) 上一个拟凝聚层同态序列正合的充要条件是其在 \(S'\) 上的逆像正合。特别地，\(S\) 上一个拟凝聚层同态是单态射、满态射或同构的充要条件，是其在 \(S'\) 上的逆像具有相应性质。限制到 \(S\) 的任意开子集上方时，此性质仍成立，并以此形式刻画忠实平坦态射。\begin{definition}{1.2}\setnumber{1.2}\label{fga3.i-b.1-definition-1.2}\skipsinglepar\par{}若 \(S\) 的每个拟紧开子集 \(U\) 的逆像拟紧，即为仿射开子集的 \emph{有限} 并，则称态射 \(\alpha \colon  S'\to  S\) 为 \emph{拟紧} 的。\end{definition}\par{}显然，只需对 \(S\) 的 \emph{仿射} 开子集验证此性质。例如，仿射态射，即仿射开子集的逆像仍仿射的态射，是拟紧的。\par{}平坦态射、忠实平坦态射以及拟紧态射这三类态射，各自对合成与“基底扩张”稳定，当然也都包含一切同构。\begin{theorem}{1}\setnumber{1}\label{fga3.i-b.1-theorem-1}\skipsinglepar\par{}设 \(\alpha \colon  S'\to  S\) 为 \emph{忠实平坦} 且 \emph{拟紧} 的预概形态射。则 \(\alpha \) 是拟凝聚层的纤维范畴 \(\mathcal {F}\)（参见 \hyperref[fga3.i-a.1-example-2]{§A, 例 2}）的一个 \emph{严格下降态射}（参见 \hyperref[fga3.i-a.1-definition-1.7]{定义~\ref*{fga3.i-a.1-definition-1.7}}）。\end{theorem}\par{}此陈述包含两件事：\begin{enumerate}
    
 \item[i.]若 \(\mathcal {F}\) 与 \(\mathscr {G}\) 为 \(S\) 上的拟凝聚层，\(\mathcal {F}'\) 与 \(\mathscr {G}'\) 为它们在 \(S'\) 上的逆像，则自然同态 \[         \operatorname {Hom}(\mathcal {F},\mathscr {G}) \to  \operatorname {Hom}(\mathcal {F}',\mathscr {G}')       \] 是从左边到右边的如下子群的双射：它由与这些层上的典范下降资料相容的同态 \(\mathcal {F}'\to \mathscr {G}'\) 组成，即沿 \(S''=S'\times _S S'\) 到 \(S'\) 的两个投影取逆像后，得到相同的同态 \(\mathcal {F}''\to \mathscr {G}''\)；

    
 \item[ii.]\(S'\) 上每个赋有关于态射 \(\alpha \colon  S'\to  S\) 的下降资料（参见 \hyperref[fga3.i-a.1-definition-1.6]{定义~\ref*{fga3.i-a.1-definition-1.6}}）的拟凝聚层 \(\mathcal {F}'\)，连同此资料，同构于 \(S\) 上的某个拟凝聚层 \(\mathcal {F}\) 的逆像。

  \end{enumerate} 在（i）中置 \(\mathcal {F}=\mathscr {O}_S\)，便得到：\begin{corollary}{1}\setnumber{1}\label{fga3.i-b.1-corollary-1}\skipsinglepar\par{}设 \(\mathscr {G}\) 为 \(S\) 上的拟凝聚层，分别以 \(\mathscr {G}'\) 和 \(\mathscr {G}''\) 表示其在 \(S'\) 和 \(S''=S'\times _S S'\) 上的逆像，\(p_1\) 和 \(p_2\) 为 \(S''\) 到 \(S'\) 的两个投影。则集合映射图式 \[       \Gamma (\mathscr {G}) \xrightarrow {\alpha ^*}       \Gamma (\mathscr {G}') \overset {p_1^*}{\underset {p_2^*}{\rightrightarrows }}       \Gamma (\mathscr {G}'')     \] 为 \emph{正合} 的（参见 \hyperref[fga3.i-a.1-definition-1.1]{定义~\ref*{fga3.i-a.1-definition-1.1}}）。\end{corollary}\par{}此外，将（i）、（ii）与 \hyperref[fga3.i-a.1-definition-1.1]{定义~\ref*{fga3.i-a.1-definition-1.1}} 结合，得到：\begin{corollary}{2}\setnumber{2}\label{fga3.i-b.1-corollary-2}\skipsinglepar\par{}设 \(\mathscr {G}\) 如 \hyperref[fga3.i-b.1-corollary-1]{推论~\ref*{fga3.i-b.1-corollary-1}} 所述。则 \(\mathscr {G}\) 的拟凝聚子层与 \(\mathscr {G}'\) 中沿两个投影 \(p_1\)、\(p_2\) 到 \(S''\) 的逆像给出 \(\mathscr {G}''\) 的同一子层的拟凝聚子层之间，存在双射对应。\end{corollary}\par{}当然，存在用商层表述的等价陈述。正如 \cref{fga3.i-a.4.e} 所见，应将 \hyperref[fga3.i-b.1-theorem-1]{定理~\ref*{fga3.i-b.1-theorem-1}} 视为 Hilbert“定理 90”的推广，它包含各种用 \(1\) 次上同调表述的形式为特例。为证明它，容易归约到 \(S=\operatorname {Spec}(A)\)、\(S'=\operatorname {Spec}(A')\) 的情形；对（i），又容易只需证明 \hyperref[fga3.i-b.1-corollary-1]{推论~\ref*{fga3.i-b.1-corollary-1}}，即对每个 \(A\) 模 \(M\)，图式 \[     M = A\otimes _A M \to      A'\otimes _A M \rightrightarrows      A'\otimes _A A'\otimes _A M   \] 正合。这由下列更一般的引理得出：\begin{lemma}{1.1}\setnumber{1.1}\label{fga3.i-b.1-lemma-1.1}\skipsinglepar\par{}设 \(A'\) 为忠实平坦 \(A\) 代数。则对每个 \(A\) 模 \(M\)，由 \(M\) 增广的复形 \(C^\bullet (A'/A,\operatorname {G_a})\otimes _A M\)（参见 \cref{fga3.i-a.4.e}）是 \(M\) 的一个 \emph{消解}。\begin{proof}\par{}只需证明，将上面的增广复形沿 \(A\) 到 \(A'\) 作基底扩张后，所得复形满足同样结论。这就归结为把 \(A\) 换成 \(A'\)，把 \(A'\) 换成 \(A'\otimes _A A'\) 后证明此陈述，因此可只考虑存在 \(A\) 代数同态 \(A'\to  A\) 的情形；几何地说，就是 \(S'\) 在 \(S\) 上有截面的情形。此时断言由 \cref{fga3.i-a.4.a} 的一般讨论得出。\end{proof}\end{lemma}\par{}顺便指出下列推论，它推广了 Galois 上同调中的熟知陈述，可与 \cref{fga3.i-a.4.e} 比较：\begin{corollary}{}\label{fga3.i-b.1-lemma-1.1-corollary}\skipsinglepar\par{}若 \(A'\) 在 \(A\) 上忠实平坦，则 \(\operatorname {H}^0(A'/A,\operatorname {G_a})=A\)，并且对 \(i\geqslant 1\)，有 \(\operatorname {H}^i(A'/A,\operatorname {G_a})=0\)。\end{corollary}\par{}为证明 \hyperref[fga3.i-b.1-theorem-1]{定理~\ref*{fga3.i-b.1-theorem-1}} 的（ii），与（i）一样，先归约到 \(S'\) 在 \(S\) 上有截面的情形，此时结果由（i）得出，参见 \cref{fga3.i-a.1.c}。\par{}显然，可以在所考虑的拟凝聚层或层系上引入各种附加结构，\emph{任意地} 改变 \hyperref[fga3.i-b.1-theorem-1]{定理~\ref*{fga3.i-b.1-theorem-1}} 及其推论。例如，在 \(S\) 上给出一个拟凝聚交换代数层，“等价于”在 \(S'\) 上给出一个这种层，并赋有关于 \(\alpha \colon  S'\to  S\) 的下降资料。考虑到 \(S\) 上此类拟凝聚层与 \(S\) 上仿射预概形之间的函子性对应，就得到下列定理的第二个断言：\begin{theorem}{2}\setnumber{2}\label{fga3.i-b.1-theorem-2}\skipsinglepar\par{}设 \(\alpha \colon  S'\to  S\) 如 \hyperref[fga3.i-b.1-theorem-1]{定理~\ref*{fga3.i-b.1-theorem-1}} 所述。则 \(\alpha \) 是一个一般不严格的 \emph{下降态射}（参见 \hyperref[fga3.i-a.2-definition-2.4]{§A, 定义 2.4}），且对预概形上的仿射概形的纤维范畴，它是 \emph{严格下降态射}（参见 \hyperref[fga3.i-a.1-definition-1.7]{§A, 定义 1.7}）。\end{theorem}\par{}定理的第一个断言意味着：设 \(X\)、\(Y\) 为 \(S\) 上的预概形，\(X'\)、\(Y'\) 为它们在 \(S'\) 上的逆像，\(X''\)、\(Y''\) 为它们在 \(S''=S'\times _S S'\) 上的逆像；则自然映射图式 \[     \operatorname {Hom}_S(X,Y) \xrightarrow {\alpha ^*}     \operatorname {Hom}_{S'}(X',Y') \overset {p_1^*}{\underset {p_2^*}{\rightrightarrows }}     \operatorname {Hom}_{S''}(X'',Y'')   \] 为 \emph{正合} 的。也就是说，\(\alpha ^*\) 给出从 \(\operatorname {Hom}_S(X,Y)\) 到 \(\operatorname {Hom}_{S'}(X',Y')\) 中与 \(X'\) 和 \(Y'\) 上的典范下降资料相容的同态所成子集的双射；即这些同态沿 \(S''\) 到 \(S'\) 的两个投影取逆像后相同。若 \(Y\) 在 \(S\) 上仿射，则这容易由 \hyperref[fga3.i-b.1-theorem-1]{定理~\ref*{fga3.i-b.1-theorem-1}} 与 \hyperref[fga3.i-b.1-corollary-1]{推论~\ref*{fga3.i-b.1-corollary-1}} 得出；一般情形需将 \hyperref[fga3.i-b.1-theorem-1]{定理~\ref*{fga3.i-b.1-theorem-1}} 与下列结果结合：\begin{lemma}{1.2}\setnumber{1.2}\label{fga3.i-b.1-lemma-1.2}\skipsinglepar\par{}设 \(\alpha \colon  S'\to  S\) 为忠实平坦拟紧态射。则 \(S\) 可等同于 \emph{\(S'\) 的拓扑商空间}；也就是说，\(S\) 的每个满足 \(\alpha ^{-1}(U)\) 为开集的子集 \(U\)，本身也是开集。\end{lemma}\par{}为补充 \hyperref[fga3.i-b.1-theorem-2]{定理~\ref*{fga3.i-b.1-theorem-2}}，还须给出 \(S'\) 预概形 \(X'\) 上下降资料的有效性判别法，此时不假设 \(X'\) 在 \(S'\) 上仿射。首先指出，\emph{这样的下降资料未必有效}；即使 \(S\) 为域 \(k\) 的谱，\(S'\) 为 \(k\) 的二次扩张 \(k'\) 的谱，而 \(X'\) 为 \(S'\) 上 \(2\) 维固有代数概形，也可能如此；Serre 利用 Nagata 的非射影曲面说明了这一点。\emph{\(X'/S'\) 上关于忠实平坦拟紧态射 \(\alpha \colon  S'\to  S\) 的下降资料有效的充要条件是，\(X'\) 是一些在 \(S'\) 上仿射且在 \(X'\) 上的下降资料下“稳定”的开子集 \(X'_i\) 的并} 若态射 \(\alpha \colon  S'\to  S\) 是 \emph{根式} 的，即单射且剩余扩张为根式扩张，则对任意 \(X'/S'\) 及 \(X'\) 上的任意下降资料，这当然成立。还可证明，若 \(\alpha \colon  S'\to  S\) \emph{有限}，并且 \(X'\) 的每个包含于 \(X'\) 在 \(S\) 上的某个纤维的有限子集，也包含于 \(X'\) 的某个在 \(S\) 上仿射的开子集，则同样成立；这就是 \emph{Weil 判别法}。特别地，当 \(X'/S'\) 拟射影时如此；此时可证明“下降所得”的预概形 \(X/S\) 也拟射影，并且若 \(X'/S'\) 射影，它也射影。归纳如下：\begin{theorem}{3}\setnumber{3}\label{fga3.i-b.1-theorem-3}\skipsinglepar\par{}设 \(\alpha \colon  S'\to  S\) 为预概形的忠实平坦拟紧态射。若 \(\alpha \) \emph{根式}，则它是 \emph{严格下降态射}。若 \(\alpha \) 有限，则相对于预概形上的拟射影、或射影预概形的纤维范畴，它是严格下降态射。\end{theorem}\begin{remark}{}\label{fga3.i-b.1-remark}\skipsinglepar\par{}我不知道在上述第二个断言中，\(\alpha \) \emph{有限} 的假设是否确实必要；可以证明，无论如何，可以“形式地”将其替换为下列看起来较弱的假设：\emph{\(S\) 的每个点都有一个开邻域 \(U\)，并存在 \(U\) 上有限且忠实平坦的 \(U'\)，以及从 \(U'\) 到 \(S'\) 的 \(S\) 态射}。上述结果未涵盖的一类情形是 \(S=\operatorname {Spec}(A)\)、\(S'=\operatorname {Spec}(\overline {A})\)，其中 \(A\) 为 Noether 局部环，\(\overline {A}\) 为其完备化；又如 \(S'\) 在 \(S\) 上拟有限，即局部同构于某个有限 \(S\) 概形的开子集，但并不有限。在这两种情形下，报告人也不知道如下问题的答案：设 \(X\) 为 \(S\) 概形，使得 \(X'=X\times _S S'\) 在 \(S'\) 上射影，是否 \(X\) 也在 \(S\) 上射影？\par{}\emph{［编者补注］} 一个拟有限、étale 且满射的态射 \(S'\to  S\)，或形如 \(\operatorname {Spec}(\overline {A})\to \operatorname {Spec}(A)\) 的态射，并不总是严格下降态射；即使 \(A\) 为代数闭域 \(k\) 上代数曲线的局部环且 \(S=\operatorname {Spec}(A)\)，也可能如此。因此，可以找到固有光滑态射 \(f\colon  X\to  S\)，使 \(X\) 成为 \(S\) 上的主 \(E\) 丛，其中 \(E\) 为椭圆曲线，并使 \(f'\colon  X'\to  S'\) 射影，而 \(f\) 不射影。这也给出 abelian 概形下的一个 \emph{非等平凡} 主齐性丛的例子。\end{remark}\section{对态射某些性质的下降的应用}\label{fga3.i-b.2}\par{}设 \(P\) 为预概形态射的一个类，\(\alpha \colon  S'\to  S\) 为预概形态射，\(f\colon  X\to  Y\) 为 \(S\) 预概形的态射，\(f'\colon  X'\to  Y'\) 为 \(f\) 沿 \(\alpha \) 的逆像。可以问，关系“\(f'\in  P\)”是否蕴含“\(f\in  P\)”？若假设 \(\alpha \) \emph{忠实平坦} 且 \emph{拟紧}，则许多重要情形下答案为肯定；后一个假设有时过强。若 \(P\) 为满射态射类或根式态射类，容易由 \(\alpha \) 的满射性直接看出；若为平坦态射类、忠实平坦态射类或光滑态射类，则由 \(\alpha \) 的忠实平坦性得出；有限型态射类的情形也容易直接看出。利用 \hyperref[fga3.i-b.1-theorem-1]{定理~\ref*{fga3.i-b.1-theorem-1}}、\hyperref[fga3.i-b.1-theorem-2]{定理~\ref*{fga3.i-b.1-theorem-2}} 和 \hyperref[fga3.i-b.1-lemma-1.2]{引理~\ref*{fga3.i-b.1-lemma-1.2}}，可见当 \(P\) 为以下各类之一时也成立：同构、开浸入、闭浸入、浸入（假设 \(f\) 有限型且 \(Y\) 局部 Noether）、仿射态射、有限态射、拟有限态射、开态射、闭态射、同胚、分离态射或固有态射。这里未涵盖的唯一重要情形是射影或拟射影态射，这已在 \cref{fga3.i-b.1} 的注中讨论。\section{沿有限忠实平坦态射的下降}\label{fga3.i-b.3}\par{}设 \(\alpha \colon  S'\to  S\) 为 \emph{有限} 态射，对应于 \(S\) 上作为模层 \emph{局部自由} 且有限型、处处非零的代数层 \(\mathscr {A}'\)。于是 \(\alpha \) 为忠实平坦拟紧态射，可以应用上述结果。在 \(S'\) 上给出拟凝聚层 \(\mathcal {F}'\)，等价于在 \(S\) 上给出拟凝聚层 \(\alpha _*(\mathcal {F}')\) 及其 \(\mathscr {A}'\) 模结构；注意 \(\mathscr {A}'=\alpha _*(\mathscr {O}_{S'})\)。为简便起见，也将 \(S\) 上的此层记为 \(\mathcal {F}'\)。类似地，\(\mathcal {F}'\) 在 \(S'\times _S S'\) 上的两个逆像 \(p_i^*(\mathcal {F}')\)，对应于拟凝聚 \((\mathscr {A}'\otimes _{\mathscr {O}_S}\mathscr {A}')\) 模层 \(\mathcal {F}'\otimes _{\mathscr {O}_S}\mathscr {A}'\) 和 \(\mathscr {A}'\otimes _{\mathscr {O}_S}\mathcal {F}'\)。给出从前者到后者的 \((S'\times _S S')\) 同态，等价于给出一个 \((\mathscr {A}'\otimes \mathscr {A}')\) 模同态；又因 \(\mathscr {A}'\) 局部自由，这等价于给出如下 \((\mathscr {A}'\otimes \mathscr {A}')\) 模同态：\[     \mathscr {U}     = \mathscr {H}\kern  -2.5pt\mathit {om}_{\mathscr {O}_S}(\mathscr {A}',\mathscr {A}')     = \mathscr {A}'\otimes \check {\mathscr {A}}'     \to  \mathscr {H}\kern  -2.5pt\mathit {om}_{\mathscr {O}_S}(\mathcal {F}',\mathcal {F}')   \] 也就是说，对 \(\mathscr {U}\) 在开子集 \(V\) 上的每个截面 \(\xi \)，给出一个 \(\mathscr {O}_S\) 模同态 \(T_\xi \colon \mathcal {F}'|V\to \mathcal {F}'|V\)，并满足条件 \begin{eqn}{3.1}\setnumber{3.1}\label{fga3.i-b.3-equation-3.1}\skipsinglepar\[       \begin {aligned}         T_{f\xi }(x)         &= fT_\xi (x),       \\T_{\xi  f}(x)         &= T_\xi (fx),       \end {aligned}     \tag{3.1}     \]\end{eqn}，其中 \(f\) 和 \(x\) 分别为 \(\mathscr {A}'\) 与 \(\mathcal {F}'\) 在包含于 \(V\) 的 \(S\) 的某个开子集上的截面。下降资料的条件（i）与（ii）（参见 \cref{fga3.i-a.1.c}）于是分别写为 \begin{eqn}{3.2}\setnumber{3.2}\label{fga3.i-b.3-equation-3.2}\skipsinglepar\[       T_{1_U}(x)       = x       \qquad \text {即 }T_{1_U}=\operatorname {id}_{\mathcal {F}'}     \tag{3.2}     \]\end{eqn} \begin{eqn}{3.3}\setnumber{3.3}\label{fga3.i-b.3-equation-3.3}\skipsinglepar\[       T_{\xi \eta }       = T_\xi  T_\eta .     \tag{3.3}     \]\end{eqn}。换言之，\emph{\(\mathcal {F}'\) 上的下降资料等价于 \(\mathscr {O}_S\) 代数层 \(\mathscr {U}=\mathscr {H}\kern  -2.5pt\mathit {om}_{\mathscr {O}_S}(\mathscr {A}',\mathscr {A}')\) 在 \(\mathscr {O}_S\) 代数层 \(\mathscr {H}\kern  -2.5pt\mathit {om}_{\mathscr {O}_S}(\mathcal {F}',\mathcal {F}')\) 中的一个表示，并满足两个线性条件（\hyperref[fga3.i-b.3-equation-3.1]{式~\ref*{fga3.i-b.3-equation-3.1}}）}。若给定 \(S'\) 上拟凝聚层的一个配对 \[     \mathcal {F}'_1\times \mathcal {F}'_2 \to  \mathcal {F}'_3   \]，可视为 \(S\) 上 \(\mathscr {A}'\) 模层的配对，以及各 \(\mathcal {F}'_i\) 上由同态 \(T_i\colon \mathscr {U}\to \mathscr {H}\kern  -2.5pt\mathit {om}_{\mathscr {O}_S}(\mathcal {F}'_i,\mathcal {F}'_i)\)（\(i=1,2,3\)）定义的粘合资料，则这些资料在显然的意义下 \emph{与给定配对相容}，当且仅当满足下列条件：\par{}对 \(\mathscr {U}\) 在某个开子集上的每个截面 \(\xi \)，以 \(\Delta \xi =\sum \xi '_i\otimes _{\mathscr {A}'}\xi ''_i\) 表示 \(\mathscr {U}\otimes _{\mathscr {A}'}\mathscr {U}\) 的如下截面：这里把 \(\mathscr {U}\) 按其左结构视为 \(\mathscr {A}'\) 模，并以公式 \[     \xi \cdot (fg)     = \sum _i\xi '_i(f)\xi ''_i(g)   \] 定义该截面，其中 \(f\)、\(g\) 为 \(\mathscr {A}'\) 在较小开子集上的截面。那么，对 \(\mathscr {A}'\) 在较小子集上的每对截面 \((x,y)\)，有公式 \begin{eqn}{3.4}\setnumber{3.4}\label{fga3.i-b.3-equation-3.4}\skipsinglepar\[       T_\xi ^{(3)}(x\cdot  y)       = \sum _i T_{\xi '_i}^{(1)}x\cdot  T_{\xi ''_i}^{(2)}y     \tag{3.4}     \]\end{eqn}。此性质也可表述为：同态 \(T^{(i)}\) 相对于给定配对 \emph{与 \(\mathscr {U}\) 的对角映射相容}。特别地，\hyperref[fga3.i-b.3-equation-3.1]{式~\ref*{fga3.i-b.3-equation-3.1}}、\hyperref[fga3.i-b.3-equation-3.2]{式~\ref*{fga3.i-b.3-equation-3.2}}、\hyperref[fga3.i-b.3-equation-3.3]{式~\ref*{fga3.i-b.3-equation-3.3}} 与 \hyperref[fga3.i-b.3-equation-3.4]{式~\ref*{fga3.i-b.3-equation-3.4}} 使我们可以用带对角映射的代数的表示来理解 \(S'\) 上拟凝聚 \emph{代数} 层上的下降资料，因此限制到交换代数，也可理解仿射 \(S'\) 概形上的下降资料。\par{}由此可得，任意 \(S'\) 预概形 \(X'\) 上的下降资料也有类似解释：给出这样的 \(X'\)，等价于给出 \emph{在 \(S\) 上} 的预概形 \(X'\)，并赋有一个 \(\mathscr {O}_S\) 代数同态 \[     \mathscr {A}'\to \mathscr {O}_{X'},   \]；而 \(X'\) 上的下降资料，等价于一个与态射 \(h\colon  X'\to  S'\) 相容的层同态 \[     \mathscr {U}     \to  \mathscr {H}\kern  -2.5pt\mathit {om}_{h^{-1}(\mathscr {O}_S)}(\mathscr {O}_{X'},\mathscr {O}_{X'})   \]，并满足类似于上面 \hyperref[fga3.i-b.3-equation-3.1]{式~\ref*{fga3.i-b.3-equation-3.1}}、\hyperref[fga3.i-b.3-equation-3.2]{式~\ref*{fga3.i-b.3-equation-3.2}}、\hyperref[fga3.i-b.3-equation-3.3]{式~\ref*{fga3.i-b.3-equation-3.3}} 和 \hyperref[fga3.i-b.3-equation-3.4]{式~\ref*{fga3.i-b.3-equation-3.4}} 的条件。\begin{example}[{Weil 的例子}]{1}\setnumber{1}\label{fga3.i-b.3-example-1}\skipsinglepar\par{}假设 \(S'/S\) 为 Galois 群为 \(\Gamma \) 的一个 \emph{Galois étale 覆盖}（参见 \cref{fga3.i-a.3} 与 \cref{fga3.i-a.4}）。则 \(S'\) 上的拟凝聚层 \(\mathcal {F}'\)，相应地 \(S'\) 预概形 \(X'\)，上的下降资料，等价于一个将 \(\Gamma \) 表示为 \((S',\mathcal {F}')\)，相应地 \((S',X')\)，的自同构的表示，并与 \(\Gamma \) 在 \(S'\) 上的作用相容。此结果是“形式的”，即可以用范畴语言证明；但从本节的观点，还得到 \(\mathscr {U}\) 的具体结构，包括其环结构、环同态 \(\mathscr {A}'\to \mathscr {U}\) 及对角映射。下列结果完全确定了此结构：\emph{作为左 \(A'\) 模，\(\mathscr {U}\) 有一组基，由对应于 \(\Gamma \) 的元素的 \(\mathscr {U}\) 的截面给出}。\end{example}\begin{example}[{Cartier 的例子}]{2}\setnumber{2}\label{fga3.i-b.3-example-2}\skipsinglepar\par{}设 \(p\) 为素数，假设 \(p\mathscr {O}_S=0\)，即 \(\mathscr {O}_S\) 的 \emph{特征为 \(p\)}；假设 \((\mathscr {A}')^p\subset \mathscr {O}_S=\mathscr {A}\)，即 \(S'/S\) \emph{为高度 \(1\) 的根式态射}；并假设 \(\mathscr {A}\) 上的代数层 \(\mathscr {A}'\) \emph{局部有一个 \(p\) 基}，即一个截面族 \((x_i)\)，使 \(\mathscr {A}'\) 作为代数由各 \(x_i\) 生成，且仅受条件 \(x_i^p=0\) 约束。假设指标 \(i\) 的集合有限，基数为 \(n\)。令 \(\mathfrak {C}\) 为 \(A'\) 的 \(A\) 导子层；它是秩为 \(n\) 的局部自由 \(A'\) 模层，此外还是 \(\mathscr {A}\) 上而非 \(\mathscr {A}'\) 上的 Lie \(p\) 代数层，并满足下列条件：\begin{eqn}{3.5}\setnumber{3.5}\label{fga3.i-b.3-equation-3.5}\skipsinglepar\[         [X,fY] = X(f)Y + f[X,Y].       \tag{3.5}       \]\end{eqn}\end{example}\begin{lemma}{}\label{fga3.i-b.3-lemma}\skipsinglepar\par{}作为赋有代数同态 \(\mathscr {A}'\to \mathscr {U}\) 的 \(\mathscr {O}_S\) 代数，\(\mathscr {U}=\mathscr {H}\kern  -2.5pt\mathit {om}_{\mathscr {O}_S}(\mathscr {A}',\mathscr {A}')\) 由左 \(A'\) 子模 \(\mathfrak {C}\) 生成，并满足下列附加关系：\begin{eqn}{3.6}\setnumber{3.6}\label{fga3.i-b.3-equation-3.6}\skipsinglepar\[         \begin {cases}           Xf-fX &= X(f)         \\XY-YX &= [X,Y]         \\X^p &= X^{(p)}.         \end {cases}       \tag{3.6}       \]\end{eqn}\end{lemma}\par{}由上述引理，\(S'\) 上的拟凝聚层 \(\mathcal {F}'\) 的下降资料，等价于对每个 \(X\in \mathfrak {C}\) 给出 \(\mathcal {F}'\) 的一个 \(\mathscr {O}_S\) 自同态 \(\overline {X}\)，满足以下条件：\begin{eqn}{3.7}\setnumber{3.7}\label{fga3.i-b.3-equation-3.7}\skipsinglepar\[       \overline {fX} = f\overline {X}     \tag{3.7}     \]\end{eqn} \begin{eqn}{3.8}\setnumber{3.8}\label{fga3.i-b.3-equation-3.8}\skipsinglepar\[       \overline {X}(fx) = X(f)x + f\overline {X}(x)     \tag{3.8}     \]\end{eqn} \begin{eqn}{3.9}\setnumber{3.9}\label{fga3.i-b.3-equation-3.9}\skipsinglepar\[       \overline {[X,Y]} = [\overline {X},\overline {Y}]     \tag{3.9}     \]\end{eqn} \begin{eqn}{3.10}\setnumber{3.10}\label{fga3.i-b.3-equation-3.10}\skipsinglepar\[       \overline {X^{(p)}} = \overline {X}^p.     \tag{3.10}     \]\end{eqn} 这可称为一个 \emph{\(\mathcal {F}'\) 上的线性联络}，它还是 \emph{平坦} 的，并且 \emph{与 \(p\) 次幂相容}。类似地可以写出 \(S'\) 预概形 \(X'\) 上下降资料的概念，只需将 \hyperref[fga3.i-b.3-equation-3.4]{式~\ref*{fga3.i-b.3-equation-3.4}} 换为各个 \(\overline {X}\) 都是 \(\mathscr {O}_{X'}\) 的 \emph{导子} 的条件。由于态射 \(S'\to  S\) 为根式态射，\hyperref[fga3.i-b.1-theorem-3]{定理~\ref*{fga3.i-b.1-theorem-3}} 保证每个这样的下降资料均有效，因而定义一个 \(S\) 预概形 \(X\)。\par{}注意，对 \(\mathcal {F}'\) 或 \(X'\) 无须作任何平坦性、非奇异性或有限性假设。\section{对有理性判别法的应用}\label{fga3.i-b.4}\par{}设 \(X\) 为 \(S\) 预概形，使得 \(\mathscr {O}_X\) 在 \(S\) 上的直像为 \(\mathscr {O}_S\)；对任意平坦基底扩张 \(S'\to  S\)，此性质仍成立。若 \(\mathcal {F}\) 为 \(X\) 上的 \emph{可逆层}，即秩为 \(1\) 的局部自由层，则 \(\mathcal {F}\) 的自同构与 \(\mathscr {O}_S\) 的可逆截面之间存在双射对应，前者可等同于 \(\mathscr {O}_X\) 的可逆截面。现设 \(s\) 为 \(X\) 在 \(S\) 上的截面，将 \(S\) 上可逆层 \(s^*(\mathcal {F})\) 的截面定义为 \emph{\(\mathcal {F}\) 在 \(s\) 上方的截面}。由前述讨论，设 \(\mathcal {F}_i\)（\(i=1,2\)）是 \(X\) 上的可逆层，各自赋有 \(s\) 上方的一个截面；那么 \emph{若 \(\mathcal {F}_1\) 与 \(\mathcal {F}_2\) 同构，则恰好存在一个从 \(\mathcal {F}_1\) 到 \(\mathcal {F}_2\) 的同构，与所给截面相容}，即将前一个截面送到后一个截面。独立于截面 \(s\)，若 \(S\) 的每个点都有开邻域 \(U\)，使 \(\mathcal {F}_1\) 与 \(\mathcal {F}_2\) 在 \(X|U\) 上的限制同构，则称 \(X\) 上的两个可逆层 \(\mathcal {F}_1\) 和 \(\mathcal {F}_2\) \emph{等价}。于是 \emph{\(X\) 上的每个可逆层 \(\mathcal {F}\) 都等价于一个在 \(s\) 上方赋有标定截面的可逆层 \(\mathcal {F}_1\)}，取 \(\mathcal {F}_1=Fs^*(\mathcal {F})^{-1}\) 即可，\emph{并且 \(\mathcal {F}_1\) 在同构意义下确定}。换言之，\(X\) 上可逆层的 \emph{模等价} 分类，与赋有一个标定截面的可逆层的 \emph{模同构} 分类相同。\par{}这些性质在平坦基底扩张 \(\alpha \colon  S'\to  S\) 下仍成立，只需将截面 \(s\) 换为其沿 \(\alpha \) 的逆像 \(s'\)；因此结合 \hyperref[fga3.i-b.1-theorem-1]{定理~\ref*{fga3.i-b.1-theorem-1}}，得到：\par{}\emph{设预概形 \(X/S\) 如上，并有截面 \(s\)；设 \(\alpha \colon  S'\to  S\) 为忠实平坦拟紧态射，\(\mathcal {F}'\) 为 \(X'=X\times _S S'\) 上的可逆层。\(\mathcal {F}'\) 等价于 \(X\) 上某个可逆层 \(\mathcal {F}\) 在 \(X'\) 上的逆像的充要条件是，其在 \(X'\times _X X'=X\times _S(S'\times _S S')\) 上的两个逆像 \(p_1^*(\mathcal {F}')\) 和 \(p_2^*(\mathcal {F}')\) 等价。此时，\(\mathcal {F}\) 在等价意义下确定。} 此时称 \(\mathcal {F}'\) 在 \(S\) 上 \emph{有理}。\par{}当 \(\alpha \colon  S'\to  S\) 如上一节的 \hyperref[fga3.i-b.3-example-1]{例~\ref*{fga3.i-b.3-example-1}} 与 \hyperref[fga3.i-b.3-example-2]{例~\ref*{fga3.i-b.3-example-2}} 时考虑此原则，就恢复了 \emph{Weil 和 Cartier 的有理性判别法}。注意，那些作者只考虑 \(S\)、\(S'\) 为域的谱的情形；此时 \(S\) 更是局部环的谱，上述等价关系恰好就是同构关系。在第一种情形，\(\mathcal {F}'\) 在 \(S\) 上有理，当且仅当其在 \(\Gamma \) 下的像均等价于 \(\mathcal {F}'\)。为表达第二种情形的有理性判别法，更一般地考虑 \(X'/X\) 的对角态射 \(X'\to  X''=X'\times _X X'\)、相应的 \(X'\times _X X'\) 上的理想层 \(\mathscr {I}\) 及层 \(\mathscr {I}/\mathscr {I}^2\)；后者可等同于其在 \(X'\) 上的逆像 \(\Omega _{X'/X}^1\)，即 \emph{\(X'\) 相对于 \(X\) 的 \(1\) 次微分层}。由于各 \(\mathcal {F}''_i=p_i(\mathcal {F}')\)（\(i=1,2\)）在对角线上的限制同构，二者均同构于 \(\mathcal {F}'\)，也就是说 \(\mathcal {F}''_1(\mathcal {F}''_2)^{-1}=\mathcal {F}''\) 在对角线上的限制平凡，因此 \(\mathcal {F}''\) 在 \((X'',\mathscr {O}_{X''}/\mathscr {I}^2)\) 上的限制，在同构意义下，由 \[     \operatorname {H}^1(X'',\mathscr {I}/\mathscr {I}^2)     = \operatorname {H}^1(X',\Omega _{X'/X}^1).   \] 的一个确定元素 \(\xi \) 给出。此外，精确地说，有 \(\Omega _{X'/X}^1=\Omega _{S'/S}^1\otimes _{\mathscr {O}_S}\mathscr {O}_X\)，所以 \emph{若 \(\Omega _{S'/S}^1\) 在 \(S\) 上局部自由}，如 Cartier 情形那样，\emph{则 \(\xi \) 定义 \(S'\) 上 \(\operatorname {R}^1f'(\mathscr {O}_{X'})\otimes \Omega _{S'/S}^1\) 的一个截面}，称为 \emph{\(X'/S\) 上可逆层 \(\mathcal {F}'\) 的 Atiyah–Cartier 类}；\emph{此类消失，当且仅当 \(\mathcal {F}'\) 沿} \[     (X'',\mathscr {O}_{X''}/\mathscr {I}^2)     = X\times _S(S'',\mathscr {O}_{S''}/\mathscr {J}^2)   \] \emph{的两个投影到 \(X'\) 的逆像等价}，其中 \(\mathscr {J}\) 是 \(S''=S'\times _S S'\) 上由对角态射 \(S'\to  S'\times _S S'\) 定义的理想层。因此，此类的消失显然是 \(\mathcal {F}'\) 在 \(X''=X\times _S S''\) 本身上的两个逆像等价的 \emph{必要} 条件，从而也是 \(\mathcal {F}'\) 等价于 \(X\) 上某个可逆层 \(\mathcal {F}\) 的逆像的必要条件。Atiyah–Cartier 类也可理解为在 \(S'\) 上局部存在 \(\mathcal {F}'\) 的关于 \(X'/X\) 的导子的一个 \emph{联络} 的障碍；只要知道对应于 \(S'/S\) 的导子到 \(X'\) 的自然延拓的 \(\mathcal {F}'\) 的导子，这个联络就确定了。由此及上一节结果，容易得出：在上述 \hyperref[fga3.i-b.3-example-2]{例~\ref*{fga3.i-b.3-example-2}} 的情形，且 \(X/S\) 有截面时，Atiyah–Cartier 类消失也是 \(\mathcal {F}'\) 在 \(S\) 上有理的充分条件。\section{对 abelian 概形的基底限制的应用}\label{fga3.i-b.5}\par{}设 \(S\) 为预概形。\(S\) 上的 \emph{abelian 概形}，定义为 \(S\) 上的一个光滑固有概形 \(X\)，使其在各点 \(x\in  S\) 处的纤维为 \(k(x)\) 上 abelian 簇的概形。假设 \(S\) 为 Noether 且 \emph{正则} 的，即其局部环均正则，则利用 Murre \hyperref[Mur1958]{[4]} 的 \emph{连通性定理}，可以证明 \emph{\(X\) 在 \(S\) 上的每个有理截面都处处有定义}，即它是一个截面，这推广了 Weil 的经典定理；至少在前引定理目前已证实的“等特征”情形下如此。由此更一般地得到：若 \(X'\) 为 \(S\) 上的光滑概形，则从 \(X'\) 到 \(X\) 的每个有理 \(S\) 映射都处处有定义。由此得到如下结果，它推广了 Chow–Lang 的一个结果：\emph{设 \(S\) 为 Noether 正则预概形，\(K\) 表示其有理函数环}，它是若干域的直和，\emph{设 \(X\) 为 \(K\) 上的 abelian 概形；若 \(X\) 同构于形如 \(X_0\times _S\operatorname {Spec}(K)\) 的 \(K\) 概形，其中 \(X_0\) 为 \(S\) 上的 abelian 概形，则 \(X_0\) 在唯一同构的意义下确定。}\par{}利用上述唯一性结果，可见将 \(X\) 的基底限制下去的问题在 \(S\) 上是局部的，因此只需知道如何限制到 \(\operatorname {Spec}(\mathscr {O}_x)\)，其中 \(x\in  S\)。同样，若 \(S'\to  S\) 为有限型 \emph{光滑} 态射，\(Y'\) 为 \(S'\) 的有理函数环，且 \(X\otimes _K K'\) 形如 \(X'_0\times _{S'}\operatorname {Spec}(K')\)，则 \emph{\(X'_0\) 赋有关于 \(\alpha \) 的典范下降资料}。考虑到 \hyperref[fga3.i-b.1-theorem-3]{定理~\ref*{fga3.i-b.1-theorem-3}}，得到：\begin{proposition}{5.1}\setnumber{5.1}\label{fga3.i-b.5-proposition-5.1}\skipsinglepar\par{}设 \(S\) 为不可约、正则且 Noether 的预概形，有理函数域为 \(Y\)；设 \(K'\) 为 \(K\) 的一个 \emph{在 \(S\) 上非分歧} 的有限扩张，\(S'\) 为 \(S\) 在 \(K'\) 中的正规化，因此为 \(S\) 的 étale 覆盖；设 \(X\) 为 \(K\) 上的 abelian 概形，使 \(X\otimes _K K'\) 形如 \(X'_0\times _{S'}\operatorname {Spec}(K')\)，其中 \(X'_0\) 为 \(S'\) 上射影的 abelian 概形。则 \(X\) 形如 \(X_0\times _S\operatorname {Spec}(K)\)，其中 \(X_0\) 为 \(S\) 上射影的 abelian 概形。\end{proposition}\begin{remark}{}\label{fga3.i-b.5-remark}\skipsinglepar\par{}报告人不知道，能否将 \(S'\to  S\) 为满射 étale 覆盖的假设——它使我们能应用 \hyperref[fga3.i-b.1-theorem-3]{定理~\ref*{fga3.i-b.1-theorem-3}}——替换为有限型 \emph{光滑} 且 \emph{满射} 态射的假设，即使假设它为 étale 态射也不知道；也不知道，若不假设 \(X'_0\) 在 \(S'\) 上射影，命题是否仍成立，而这个条件也许自动满足。\end{remark}\section{对局部平凡性与等平凡性判别法的应用}\label{fga3.i-b.6}\par{}设 \(S\) 为预概形，\(G\) 为 \(S\) 上的一个“\emph{群预概形}”，\(P\) 为 \(S\) 上的一个预概形，其上“\emph{\(G\) 作用}”，设为右作用。若由 \(G\) 在 \(P\) 上的作用诱导的熟知态射 \[     G\times _S P \to  P\times _S P   \] 是 \emph{同构}，则称 \(P\) 关于 \(G\) 是 \emph{形式主齐性} 的。以下假设 \(G\) 在 \(S\) 上 \emph{平坦}，从而在 \(S\) 上忠实平坦；并只将 \(S\) 上 \emph{忠实平坦} 且 \emph{拟紧} 的形式主齐性丛 \(P\) 称为 \(G\) 的 \emph{主齐性丛}。显然，这等价于能找到基底 \(S\) 的一个 \emph{忠实平坦} 且 \emph{拟紧扩张} \(S'\to  S\)，使 \(G'=G\times _S S'\) 的形式主齐性丛 \(P'=P\times _S S'\) \emph{平凡}，即同构于 \(G'\)，也就是有一个截面；特别地，可以取 \(S'=P\)。还要指出，若 \(S\) 局部 Noether，则对 \(P\) 的忠实平坦性假设，等价于对每个 \(s\in  S\)，\(\overline {P}_s=P\times _S\operatorname {Spec}(\overline {\mathscr {O}}_s)\) 在 \(\overline {\mathscr {O}}_s\) 上忠实平坦；其中 \(\overline {\mathscr {O}}_s\) 表示局部环 \(\mathscr {O}_s\) 的完备化，这由 \(\overline {\mathscr {O}}_s\) 在 \(\mathscr {O}_s\) 上忠实平坦这一事实得出。此外，若 \(P\) 在 \(S\) 上有限型且 \(S\) 局部 Noether，则满足上述条件的点 \(s\) 的集合可构造；所以若 \(S\) 为“Jacobson 预概形”，例如域上，或更一般地 Jacobson 环上的有限型概形，则只需对 \(S\) 的 \emph{闭} 点验证所述条件。这使我们归结到基底为完备局部环 \(A\) 的谱的情形。若 \(S=\operatorname {Spec}(A)\)，其中 \(A\) 为完备 Noether 局部环，且 \(P\) 在 \(S\) 上有限型，则 \(P/S\) 的忠实平坦性还等价于存在一个在 \(S\) 上 \emph{有限且平坦} 的 \(S'\)，使 \(P'\) 平凡；若此外 \(G\) 在 \(S\) 上 \emph{光滑}，则可假设 \(S'\) 在 \(S\) 上 \emph{étale}。此时，若再假设 \(A\) 的剩余域代数闭，即“\emph{几何情形}”，则 \(P\) 在 \(A\) 上忠实平坦，当且仅当它平凡。因此，若 \(S\) 为代数闭域上的代数预概形，且 \(G\) 在 \(S\) 上光滑且有限型，则 \(S\) 上的忠实平坦性条件等价于 Serre［\hyperref[Ser1958a]{6}，第 1—12 页］的解析平凡性条件（SLF）。\par{}还可以对 \(P\) 考虑其他更强的、具有“局部平凡性”性质的条件。特别地，若对每个 \(s\in  S\)，存在 \(s\) 的一个开邻域 \(U\)，以及一个 \emph{有限且忠实平坦} 态射，相应地一个 \emph{满射 étale 覆盖}，\(U'\to  U\)，使 \(P'=P\times _S U'\) 平凡，则称 \(P\) 为 \emph{等平凡} 的，相应地为 \emph{严格等平凡} 的。这里与 Serre \hyperref[GD1960]{[1]}\footnote{译校注：原刊此处写作“Serre [1]”，但本报告原书目 [1] 为 Dieudonné–Grothendieck 的 EGA；疑指 Serre 的 [6]。保留原编号，不据猜测改换引文。} 的术语不同，他用“局部等平凡”指我们所谓的“严格等平凡”。严格等平凡性主要在 \(G\) 在 \(S\) 上光滑时有用，但在其他情形是不适当的概念。\par{}若 \(G\) 在 \(S\) 上 \emph{仿射}，则由 \cref{fga3.i-b.2}，\(G\) 的每个主齐性丛 \(P\) 均仿射；因此借助 \hyperref[fga3.i-b.1-theorem-2]{定理~\ref*{fga3.i-b.1-theorem-2}}，能够沿忠实平坦拟紧态射将这种丛“下降”。特别地，取 \(G=\operatorname {GL}(n)_S\)，它由如下条件定义：\(S\) 预概形上的取值于群范畴的函子 \(S'\mapsto \operatorname {Hom}(S',G)\) 可等同于 \cref{fga3.i-a.4} 所述函子 \(\operatorname {GL}(n)(S')=\operatorname {GL}(n,\operatorname {H}^0(S',\mathscr {O}_{S'}))\)。利用以下事实：\begin{enumerate}
    
 \item[i.]\(G\) 的每个主齐性丛，相应地 \(S\) 上的每个秩为 \(n\) 的局部自由层，在适当的 \(S\) 的忠实平坦拟紧扩张后，都同构于“平凡”对象 \(G\)，相应地 \(\mathscr {O}_S^n\)；


    
 \item[ii.]可以沿这种态射下降所考虑类型的对象，即 \(G\) 的主齐性丛，相应地秩为 \(n\) 的局部自由层；最后，


    
 \item[iii.]任意 \(S'/S\) 上的平凡丛的自同构群函子性地同构于 \(S'\) 上秩为 \(n\) 的平凡局部自由层的自同构群，

  \end{enumerate} 可以“形式地”得出，在 \(S\) 上，或在某个 \(S'/S\) 上，给出群 \(G\) 的一个主齐性丛，“等价于”给出一个秩为 \(n\) 的局部自由层。更精确地说，有一个 \emph{纤维范畴的等价}。特别地，得到：\begin{proposition}{6.1}\setnumber{6.1}\label{fga3.i-b.6-proposition-6.1}\skipsinglepar\par{}群 \(\operatorname {GL}(n)_S\) 的每个主齐性丛都局部平凡。\end{proposition}\par{}由熟知的论证，对其他结构群，例如 \(\operatorname {SL}(n)_S\)、\(\operatorname {Sp}(n)_S\) 及这些群的乘积，也得到相同结果。因此也得到：若 \(F\) 为 \(G=\operatorname {GL}(n)_S\) 的一个在 \(S\) 上平坦的闭子群，使商 \(G/F\) 存在，且 \(G\) 是 \(G/F\) 上结构群为 \(F\times _S(G/F)\) 的等平凡、相应地严格等平凡的主齐性丛，则结构群为 \(F\) 的 \emph{每个} 主齐性丛都等平凡、相应地严格等平凡。这适用于迄今使用过的 \(S\) 上的所有“线性群”；特别适用于 \(G=S\times _k\Gamma \) 的情形，其中 \(S\) 为域 \(k\) 上的预概形，\(\Gamma \) 为 \(k\) 上经典意义下的线性群，特别地它光滑。这就为这种群回答了 Serre 的一个问题（\emph{同前引处}）。\par{}还要指出，对我们所知道的大多数 \(S\) 上的单纯群，无论是否线性，并且当然对上述形如 \(S\times _k\Gamma \) 的所有群，都可证明每个等平凡主齐性丛都严格等平凡。这特别回答了 Serre 的另一个问题（\emph{同前引处}，第 1—14 页）；这里考虑到，沿 Cartier \emph{式方法} 下降（参见 \hyperref[fga3.i-b.3-example-2]{例~\ref*{fga3.i-b.3-example-2}}）所得的主齐性丛显然等平凡。\begin{remark}{}\label{fga3.i-b.6-remark}\skipsinglepar\par{}这些问题的一个本质困难，撇开商概形 \(G/F\) 的存在性问题，就是缺少沿忠实平坦 \emph{非有限} 态射的下降资料的有效性判别法。\end{remark}
\section*{原报告参考文献}
\begin{enumerate}
\item[1.]\setnumber{1}\label{fga3.i-bib-1} DIEUDONNÉ (J.) et GROTHENDIECK (Alexander). \emph{Éléments de géométrie algébrique}，即将发表于 \emph{Publications mathématiques de l'Institut des Hautes Études Scientifiques}。
\item[2.]\setnumber{2}\label{fga3.i-bib-2} GRAUERT (H.) und REMMERT (R.). \emph{Komplexe Räume}，\emph{Math. Annalen}，138 (1958)，245--318。
\item[3.]\setnumber{3}\label{fga3.i-bib-3} GROTHENDIECK (Alexander). \emph{Géométrie formelle et géométrie algébrique}，\emph{Séminaire Bourbaki}，11 (1958/59)，第 182 次报告。
\item[4.]\setnumber{4}\label{fga3.i-bib-4} MURRE (J. P.). \emph{On a connectedness theorem for a birational transformation at a simple point}，\emph{Amer. J. Math.}，80 (1958)，3--15。
\item[5.]\setnumber{5}\label{fga3.i-bib-5} SERRE (Jean-Pierre). \emph{Géométrie algébrique et géométrie analytique}，\emph{Ann. Inst. Fourier Grenoble}，6 (1955--56)，1--42。
\item[6.]\setnumber{6}\label{fga3.i-bib-6} SERRE (Jean-Pierre). \emph{Espaces fibrés algébriques}，\emph{Séminaire Chevalley}，2，1958：\emph{Anneaux de Chow et applications}，第 1 次报告。
\end{enumerate}
